\documentclass[11pt,a4paper]{article}

\usepackage[T1]{fontenc}
\usepackage[utf8]{inputenc}
\usepackage{lmodern}
\usepackage[margin=2.4cm]{geometry}
\usepackage{microtype}
\usepackage{amsmath,amssymb,amsthm}
\usepackage{booktabs}
\usepackage{array}
\usepackage{longtable}
\usepackage{enumitem}
\usepackage{titlesec}
\usepackage[hidelinks]{hyperref}
\usepackage{xcolor}

\definecolor{rulegrey}{gray}{0.45}
\titleformat{\section}{\large\bfseries}{\thesection}{0.7em}{}
\titleformat{\subsection}{\normalsize\bfseries}{\thesubsection}{0.7em}{}
\titleformat{\subsubsection}{\normalsize\itshape}{\thesubsubsection}{0.7em}{}
\titlespacing*{\section}{0pt}{1.6em}{0.6em}
\titlespacing*{\subsection}{0pt}{1.2em}{0.4em}
\setlist[itemize]{leftmargin=1.2em,itemsep=0.25em,topsep=0.3em}
\setlist[enumerate]{leftmargin=1.5em,itemsep=0.35em,topsep=0.4em}
\newcommand{\hi}[1]{\textbf{#1}}
\newcommand{\sev}[1]{\textsc{#1}}
\newcommand{\Frob}[1]{\lVert #1 \rVert_{F}}
\newcommand{\Hen}[1]{\mathrm{He}_{#1}}
\newcommand{\E}{\mathbb{E}}
\newcommand{\R}{\mathbb{R}}
\newcommand{\Sg}{S^{\mathrm{G}}}
\newcommand{\tstar}{t_\ast}
\newcommand{\lmax}{\lambda_{\max}}
\newtheorem{theorem}{Theorem}
\newtheorem{proposition}[theorem]{Proposition}
\newtheorem{corollary}[theorem]{Corollary}
\newtheorem{lemma}[theorem]{Lemma}
\theoremstyle{definition}
\newtheorem{assumption}[theorem]{Assumption}
\theoremstyle{remark}
\newtheorem{remark}[theorem]{Remark}

\begin{document}

\begin{center}
{\LARGE\bfseries What survives the audit}\\[0.45em]
{\large A scientific audit of the ICML manuscript and of the 18 September report,\\ with the analyses it forced and the paper it recommends}\\[0.9em]
{\small Internal research report \quad$\cdot$\quad 18--19 September 2026}\\[0.25em]
{\footnotesize Every number that is new here comes from a script under \texttt{experiments/} or \texttt{outputs/E1x/} named in the text; every number that is old comes from a frozen run named in the text.}
\end{center}

\vspace{0.4em}
\hrule height 0.8pt
\vspace{1.2em}

\noindent
This report audits two documents as primary sources: the ICML manuscript \emph{What Survives Diffusion Noise?} (\texttt{paper/main.pdf}, 18 September) and the internal report \emph{The EM chapters, re-examined} (\texttt{reports/em-chapters-2026-09-18.pdf}). It does what the brief asked, in the order the brief asked: science first, prose last. Section~1 says what stands, what does not, and what changed; Section~2 lists the issues by severity; Sections~3--5 are the mathematical, experimental and consistency audits; Section~6 reports the analyses run for this audit (all on the laptop, all cheap, several decisive); Sections~7--13 are the forward plan---experiments, claims, narrative, structure, figures, the two referee reports, and the checklist.

\medskip
\noindent\hi{The one-paragraph version.} The theory is right and the constants are right, but Theorem~3 is stated under a moment condition its proof does not use: the proof needs an exponential moment, which every innovation family in the paper except Student-$t$ has, and Student-$t$ is precisely the family that departs. The heavy-tailed departure the manuscript attributes to the moment condition is, on the grid it was measured on, a truncation artefact: on a wide enough domain the sampled Student-$t_5$ chain reproduces the $e^{-8t}$ constant ($1.03$ $[0.99,1.09]$ at $t=4$ from $8192$ sequences), while an untruncated scalar quadrature shows the \emph{expected} gap decaying as $e^{-\nu t}$, carried by events no sample contains---so the moment condition is real, and its violation is invisible in practice. The ``all belief mass outside the domain'' diagnostic is one site of one chain in ten thousand, not a property of the configuration. The $7.3\times$ ``information-matched'' comparison gives the network the clean sequence as a regression target and is therefore not information-matched; the genuinely matched arm (implicit score matching, no clean target) gives $7.9\times$ $[6.6, 9.3]$, so the headline survives with a larger number and a different sentence. The ``estimator attains the Cram\'er--Rao bound at $1.02\times$'' claim does not survive: at the levels it averages over, the fitted shape is a mixture of a shrunken cluster below the bound and Gaussian-collapsed fits at the coordinate's range, and $1.02$ is the mean of ratios from $0.7$ to $1.5$. The capacity reversal, the family $V$, the regime sweep, generation and the misspecification brackets stand. Four new results were obtained for this audit and belong in the paper: the exact radius-$r$ locality floor, which at $\rho=0.85$ is $0.030$ for radius $4$ against the window head's $0.118$ and withdraws ``locality is a ceiling'' (the plateau is the architecture's); a closed-form asymptotic bias for every misspecification mechanism that reproduces the $M=2048$ EM--BP error to three decimals and predicts the order in $\varepsilon$ at which each mechanism bites; an exact fourth-order (Edgeworth) correction to the covariance twin that closes $96\%$ of the closure gap at $t=0.68$ and $100\%$ by $t=2$ on every family with an exponential moment, and cannot close it on Student-$t$; and the exact information curve at $\rho=0.5$, which says the sample-complexity experiment E12b as pre-registered would have declared the corollary refuted by a slope of $6.5$ that is in fact its prediction on those levels (amended before any data, Section~7).

\tableofcontents

%% ------------------------------------------------------------------------------------
\section{Executive scientific assessment}
\label{sec:exec}

\subsection{What is genuinely strong}

\begin{enumerate}
\item \hi{The laboratory.} A non-Gaussian, correlated law in $\R^{32}$ whose diffusion score is computable exactly, with a validated numerical instrument (at the frozen configuration the reference drifts $9\times10^{-6}$ under an eightfold refinement and a $1.5\times$ widening of the domain). Nothing comparable exists in the solvable-diffusion literature: the Gaussian, mixture and hierarchical-discrete models are either linear in the score or nonlinear only through a discrete latent.
\item \hi{Theorem 2 and the algebraic skeleton of Theorem 3.} Hermite attenuation is exact and elementary; the four constants of Theorem 3 ($e^{-mt}$ per degree, $\kappa_4$ as the degree-4 difference, $\tfrac12$ in the covariance rate, $\tfrac16$ in the score gap) are machine-checked in Lean, and the fifth---$\tfrac1{24}$ in the shape rate---is verified in this audit by scalar quadrature to $0.25\%$ (Section~6.5). The closure-gap constant is reproduced on the chain to $1$--$6\%$ for Laplace and bimodal innovations and, on a wide enough domain, to $1\%$ for Student-$t_5$ at $t=4$.
\item \hi{Proposition 1 as an apparatus check.} ``$C=1$ lands on the LMMSE floor'' is verified at fifteen $(\rho, L)$ cells to $0.01$ and across a one-parameter family at every member. It is the single most reassuring fact about the whole pipeline (grid, message passing, EM, scoring), and it is stated correctly.
\item \hi{The capacity reversal and its diagnosis.} Retiring ``capacity does not help'' with the exact mechanism (a stopping rule that fired for one arm) and the two controls that isolate it (E02a\_v1 at a shorter ladder, E10 in the thesis's own estimand) is exemplary. It should be foregrounded, not hidden.
\item \hi{The experimental discipline.} Frozen protocol, pre-registration, hashed bundles, paired seeds, a generator that refuses dirty runs, a manuscript check that refuses typed decimals, and an independent re-derivation of every headline number. This audit found no arithmetic error in any frozen number. Every problem found is in what a number was \emph{said} to mean.
\item \hi{Misspecification as a measured boundary.} Four mechanisms, sixteen seeds, brackets rather than fitted crossings, and the separation of estimation from bias by two sample sizes. With the closed-form bias of Section~6.3 it becomes a small theory.
\end{enumerate}

\subsection{What is currently vulnerable}

\begin{enumerate}
\item \hi{Theorem 3's hypothesis.} Stated for finite eighth moments; proved (in Appendix A.3) with an analyticity of $\Lambda(s,b)$ at the origin and a tail bound $P_t(|\alpha x|>1)\le e^{-c/\alpha^2}$ that hold only under an exponential moment. A theoretical referee will find this in ten minutes. Section~3.3 gives the repair: state the theorem under $\E e^{\lambda|a|}<\infty$; keep finite moments for Theorem~2; add the heavy-tailed regime as a proposition with its own rate.
\item \hi{``Information-matched''.} The arm behind $7.3\times$ trains by denoising score matching on $(a, x_t)$ pairs---the network sees the clean chain. Its sigma-algebra strictly contains EM--BP's. The genuinely matched arm (implicit score matching on $x_t$ alone) exists in the same frozen file, gives $7.9\times$ $[6.6, 9.3]$ paired over sixteen seeds, and is the number to report (Section~4.3).
\item \hi{``Attains the Cram\'er--Rao bound''.} The $1.02\times$ average pools levels where the estimator is unimodal and near the bound ($t\le0.4$) with levels where it is a mixture of a shrunken cluster and collapsed fits ($t=0.7$: four of sixteen at excess kurtosis $\approx0$, the other twelve with dispersion $0.76$ against a bound of $1.79$). The bound is for the one-parameter shape family, the estimator is an eight-component mixture with early stopping selected on an oracle. None of this is efficiency in the statistical sense. Section~4.5 gives the defensible statement.
\item \hi{Oracle model selection.} EM--BP, the standard-recipe networks and the fixed-corruption arm all select their checkpoint on a validation bundle scored against the \emph{true-kernel} score. Only the implicit-score-matching arm selects without an oracle. A control with held-out log-evidence selection is running (Section~6.7); at $M=2048$ the validation curve is monotone through the ladder, so the number will not move, but the sentence must say what the selection uses.
\item \hi{The Student-$t$ envelope story.} ``Essentially all of its belief mass outside the domain'' is a maximum over $8192$ (chain, site) pairs at which exactly one site (seed 2: $|a|=15.4$; seeds 0, 1, 3: one site each beyond $8$) is pinned to the boundary. The instrument finding is real and better than stated---for polynomial tails the truncation error is a Poisson number of exceedances, and no finite domain removes it---but the sentence is wrong and the E07 ladder's Student-$t$ rows are wrong by a few sites, not by ``all''.
\item \hi{Locality as a ceiling.} The claim rests on a plateau ($0.121\to0.118$) and an argument (``a radius-4 window cannot represent a message that has crossed the chain''). The exact best radius-$r$ estimator is computable (Section~6.4) and at $\rho=0.85$ its pooled error at radius $4$ is $0.030$ (peaking at $0.065$ at intermediate levels), against the window head's $0.118$ ($0.075$ on the interior sites) and EM--BP's $0.023$. Locality costs little information; the plateau is the architecture's. The sentence must change.
\item \hi{Overstatements of what E06 confirmed.} The introduction says the $e^{4t}$ separation ``is confirmed''; the contribution list says it is ``measured''; Section~4's heading asserts it. E06 saturates before the asymptotic window at $\rho=0.85$, E06b is inconclusive, E12 is confounded. What is established is the rates \emph{by computation}, the ordering \emph{by measurement}, and the collapse level. Say that.
\item \hi{Stale and inconsistent text.} A \texttt{[PENDING]} note in the abstract, a dangling \verb|\cref{thm:gap}| (``??'' on page 7), a provenance paragraph saying under-resolved fits are refused while Table~1 prints them with daggers, sixteen seeds in Section~5 against eight in the regime sweep, and a related-work paragraph whose ``what is new'' list is longer than the evidence.
\end{enumerate}

\subsection{What changed because of this audit}

\begin{itemize}
\item Theorem~3 restated under an exponential-moment assumption with a corrected tail argument; a general-order theorem replaces Remark~4; the sample-complexity corollary becomes a local asymptotic minimax statement; a heavy-tailed proposition is added (Section~3).
\item The headline matched-interface number becomes $7.9\times$ (implicit score matching), with the fixed-corruption arm demoted to ``even given the clean sequence''.
\item The efficiency paragraph is rewritten around dispersion, bias and collapse fraction per level.
\item The Student-$t$ cell is described by its exceedance count; the E07 rows are recomputed on a wide domain (E07b/E07c queued; the local check of Section~6.2 anticipates them).
\item Figure~2c is redrawn from the wide-domain computation with the large-sample constants overlaid (Section~6.2--6.3): every family reaches the constant on the sampled chain; the heavy-tailed regime lives in the expectation and is stated as such.
\item Three new results enter the paper: the closed-form misspecification bias (Section~6.3), the fourth-order correction as a moment-method baseline (Section~6.6), and the moment estimators as the high-noise competitor (Section~6.8).
\item E12b's analysis plan is amended before any of its cells has run: the target is the exact information curve, whose slope on the swept levels is $6$--$7$, not $8$ (Section~6.9).
\end{itemize}

%% ------------------------------------------------------------------------------------
\section{Critical issues}
\label{sec:issues}

Severity: \sev{blocker} = a referee can reject on it; \sev{major} = must be fixed before submission; \sev{moderate} = must be fixed, low risk; \sev{minor} = editorial. ``Exp.'' says whether a new experiment is required.

{\setlength{\tabcolsep}{3pt}\small
\begin{longtable}{@{}>{\raggedright\arraybackslash}p{1.75cm}>{\raggedright\arraybackslash}p{2.0cm}>{\raggedright\arraybackslash}p{6.05cm}>{\raggedright\arraybackslash}p{4.3cm}>{\raggedright\arraybackslash}p{0.9cm}@{}}
\toprule
sev. & where & problem and reason & fix & exp. \\
\midrule
\endhead
\sev{blocker} & Thm 3, App.\ A.3 & Hypothesis ``finite eighth moments''; proof uses $P_t(|\alpha x|>1)\le e^{-c/\alpha^2}$ (false under finite moments: by Markov the tail is $O(\alpha^{16})$, polynomial) and real-analyticity of $\Lambda(s,b)$ at $(0,0)$ (false when $\E e^{s\cdot a}=\infty$ for all $s\ne0$). The remainder bound $C\alpha^6(1+|x|)^6$ on $\{|\alpha x|\le1\}$ is asserted, not proved, and the tilted moments it needs grow like $e^{|x|^2/2\Delta}$ on that set. & Assume $\E e^{\lambda|a|}<\infty$ for some $\lambda>0$ (holds for Gaussian, Laplace, uniform, mixture, generalised Gaussian $\beta\ge1$). Then $\Lambda$ is analytic in $|s|<\lambda$, $b\ge0$, the bulk is $\{|s|\le\lambda/2\}$, and the tail bound is exponential. Student-$t$ is outside the theorem and gets Proposition~\ref{prop:heavy}. & no \\
\sev{blocker} & \S5 ``matched information'', abstract, report \S5.2 & The $7.3\times$ arm is DSM on $(a,x_t)$: the network regresses on the clean chain. Not information-matched. & Headline: implicit score matching on $x_t$ only, $7.9\times$ $[6.6,9.3]$, 16/16 seeds; the fixed-corruption arm reported as ``even with the clean target, $7.3\times$''. Macro logic in \texttt{make\_baseline\_numbers.py} picks the best arm by risk; it must pick the ISM arm by definition. & no \\
\sev{major} & \S4 efficiency & ``$1.02\times$ the bound'' averages ratios $0.7$--$1.5$ over post-hoc-selected levels; the fitted shape is bimodal from $t=0.7$ (4/16 collapsed) and collapsed at $t\ge1$ (15/16); the surviving cluster's dispersion is $0.42\times$ the bound; the bound is for a different family. & Report per level: collapse fraction, dispersion and bias of the non-collapsed cluster, the one-parameter bound as a reference. Claim: unimodal and within $20\%$ of the reference for $t\le0.4$; collapse begins at $t=0.7$ and is total by the level at which the bound crosses the range. & E06c (running) \\
\sev{major} & \S5, ``locality is a ceiling'' & The exact radius-4 floor at $\rho=0.85$ is $0.030$ pooled; the window head sits at $0.118$ ($0.075$ interior). The plateau is not the radius's. & Rewritten against the floor (E15): locality costs $0.030$, the head's remainder is the architecture and its training; Fig.~6a. & no (done) \\
\sev{major} & App.\ B, App.\ E, Fig.\ 6, report \S8.2 & ``essentially all of its belief mass outside the domain'' is a max over sites; the cell has one to two sites beyond $8$ in $8192$. The drift is those sites. & State the exceedance count and the mechanism; recompute the E07 row on a domain wide enough for the drawn sample (local check in Section~6.2; E07b/c queued). & queued \\
\sev{major} & Fig.\ 2c, \S3 verification & Student-$t_5$ ``departs as the moment condition predicts'' at $t\approx3$: on the $A=8$ grid this is truncation of the prior's tail; on $A=16$ the constant is reached ($0.99$ at $t=4$) and the departure begins at $t\approx4.5$. & Replace the curve by the wide-domain run; state the departure as the heavy-tail regime (Prop.~\ref{prop:heavy}). & no (done) \\
\sev{major} & \S1, contributions (iii), \S4 heading & ``the experiment confirms'' the $e^{4t}$ separation; it does not (saturation, E06b inconclusive, E12 confounded). & Rates by computation, ordering by measurement, collapse level measured; exponents from estimator error not claimed until E12b. & E12b (queued) \\
\sev{major} & Cor.\ 5 & ``needs $M\propto e^{4t}$'' from a Fisher information alone. & State as a local asymptotic minimax lower bound (H\'ajek--Le Cam) attained by the MLE asymptotically; ``local'' precision; regime $M\to\infty$ at fixed $t$. & no \\
\sev{major} & Remark 4 & General order is a remark without hypotheses or constant. & Theorem~\ref{thm:general} below: constant $\sum_{|\nu|=m}(\partial_\theta\E H_\nu)^2/\nu!$, proved by the same grading under the exponential moment. & no \\
\sev{moderate} & All arms except ISM & Checkpoint selection on a true-score validation bundle. & Say so in the protocol; report the held-out-evidence control for EM--BP (Section~6.7). & running \\
\sev{moderate} & \S5g & \verb|\cref{thm:gap}| unresolved (``??''). & Label. & no \\
\sev{moderate} & App.\ C vs Table 1 & ``generator refuses to emit \ldots cells whose narrowest fitted component is under two grid cells'' vs daggers in Table~1 and ``1 of 320 \ldots reported, not dropped''. & Policy: reported with a flag; the generator refuses only dirty trees and incomplete designs. & no \\
\sev{moderate} & \S3 eq.\ (5) verification & Constant of the shape rate is not reproduced on the chain ($0.81$ at $\beta=2$, $0.92$ at $\beta=3$ at the plateau); the text says ``verifies the rates and the constants''. & The scalar constant is verified to $0.25\%$ (Section~6.5); on the chain the outer-product estimator's floor limits the check to $10$--$20\%$; say which constants are verified where. & no \\
\sev{moderate} & \S5 protocol vs \S5g & Sixteen seeds stated as the unit; the regime sweep has eight; E14 has eight. & State per experiment. & no \\
\sev{moderate} & E08 text & ``the screened window head''; E08 trains a radius-2, width-64 head for 16k steps, not E03's screened radius-4 winner. & Name it. & no \\
\sev{moderate} & E12 (frozen) & Confounded at $t=1$; kurtosis non-monotone in $M$; pre-registered E12b targets exponents that the information curve says are $6$--$7$ on those levels. & Amendment recorded before data (Section~7). & queued \\
\sev{minor} & Abstract & \texttt{[PENDING]}; ``Correlation \ldots decays as $e^{-2t}$'' (amplitude) beside ``information decays as $e^{-4t}$'' without the dictionary. & Rewrite; add the amplitude/information/samples dictionary once. & no \\
\sev{minor} & \S6 related work & ``what is not in that literature'' lists five items; two are specialisations. & Claim the synthesis precisely (Section~13 of the brief; here Section~8.3). & no \\
\bottomrule
\end{longtable}}

%% ------------------------------------------------------------------------------------
\section{Mathematical audit}
\label{sec:math}

Conventions: $\alpha=e^{-t}$, $\Delta=1-\alpha^2$, $x=\alpha a+\sqrt\Delta z$, $\gamma$ the standard Gaussian on $\R^L$, $\Hen{n}$ probabilists' Hermite polynomials, $H_\nu(x)=\prod_i\Hen{\nu_i}(x_i)$, $\E_\gamma[H_\nu H_\mu]=\nu!\,\delta_{\nu\mu}$. Four objects that the manuscript sometimes conflates are kept apart throughout:
\begin{description}[leftmargin=1.6em,itemsep=0.2em]
\item[(L1)] exact functional sum-product on the chain posterior;
\item[(L2)] Gaussian-projected functional sum-product (the closure of Proposition~1);
\item[(L3)] the grid implementation of (L1) at $(N_g, A)$;
\item[(L4)] the LMMSE estimator under the true covariance.
\end{description}

\subsection{Proposition 1: Gaussian closure equals the LMMSE of the covariance twin}

\paragraph{Statement audited.} (L2) with $\rho\ne0$ returns $e^{-t}\Sigma_0\Sigma_t^{-1}x$ for every innovation law of mean $0$ and variance $q$.

\paragraph{Forward induction.} Let the projected forward message at site $i-1$ be $N(m,v)$. Multiplying by the unary $\ell_{i-1}(a)\propto\exp[-(x_{i-1}-\alpha a)^2/2\Delta]$ gives a Gaussian with precision $1/v+\alpha^2/\Delta$ and mean $(m/v+\alpha x_{i-1}/\Delta)/(1/v+\alpha^2/\Delta)$, the same in both chains. The transition integral $\int\varphi(a-\rho a')g(a')\,da'$ is the density of $\rho a'+\varepsilon$ with $a'\sim g$, $\varepsilon\sim\varphi$ independent; its mean is $\rho m'$ and its variance $\rho^2v'+q$ for \emph{every} $\varphi$ with moments $(0,q)$. Projection returns $N(\rho m',\rho^2v'+q)$, which is the Gaussian chain's message. Correct.

\paragraph{Backward induction.} The backward message as a function of $a$ is $h(\rho a)$ with $h(u)=\int\varphi(a''-u)g(a'')\,da''$, the density of $a''-\varepsilon$ evaluated at $u$; $h$ integrates to one in $u$, has mean $m'$ and variance $v'+q$. As a function of $a$ it is $h(\rho a)$: its mass is $1/|\rho|$, its normalised mean $m'/\rho$ and variance $(v'+q)/\rho^2$. This is where $\rho\ne0$ is used, and it is used correctly: at $\rho=0$ the backward message is constant and has no Gaussian projection, and the Gaussian chain's backward message is likewise flat. The terminal message $\overleftarrow m_L\equiv1$ is never projected (the induction starts from $\ell_L$); the manuscript says so.

\paragraph{Normalisation constants.} Irrelevant to the posterior mean; the projection is applied to normalised shapes. Correct.

\paragraph{Belief and the LMMSE identity.} Product of three Gaussians (forward, unary, backward) is the Gaussian chain's belief; its mean is the Gaussian chain's posterior mean, which equals $\mathrm{Cov}(a,x)\mathrm{Cov}(x)^{-1}x=\alpha\Sigma_0(\alpha^2\Sigma_0+\Delta I)^{-1}x$. Correct.

\paragraph{Stationarity.} Not assumed. The initial law is $N(0,1)$ at site~1 and the induction carries site-dependent $(m,v)$; $\Sigma_0$ is the chain's covariance whatever the start. The text's ``$\mathrm{Cov}(a_i,a_j)=\rho^{|i-j|}$'' holds because $a_1$ has unit variance and $q=1-\rho^2$; nothing in the proposition needs it.

\paragraph{Layers.} The proposition is a statement about (L2) and (L4). The manuscript then uses it in two further ways that need their own one-line justifications, both true and both currently implicit:
\begin{itemize}
\item \emph{``every Gaussian-innovation kernel is bounded by the LMMSE floor''} (App.~F): BP with any Gaussian kernel is a linear map of $x$ (the Gaussian chain's smoother with that kernel's covariance), and (L4) is the minimiser of the posterior-mean MSE over all linear maps; the score error is $(\alpha/\Delta)^2$ times the mean error. So the floor bounds every linear estimator. This is a different fact from the proposition and should be stated as a corollary.
\item \emph{``$C=1$ lands on the floor''} (Table~1, E13): needs additionally that the Gaussian quasi-MLE of $(\rho,q)$ from noised data is consistent for the true $(\rho,q)$. It is: the Gaussian likelihood depends on the data only through the sample covariance, whose expectation $\alpha^2\Sigma_0+\Delta I$ identifies $(\rho,q)$. Say so in one sentence; it explains why the identity holds to $0.01$ at $M=512$ and is what makes E13's fifteen cells an apparatus check.
\end{itemize}

\paragraph{Verdict.} \sev{verified} as stated. Add the two one-line corollaries. The Lean coverage is honest about what it checks (the moment maps for a \emph{Gaussian} innovation, not the moments-only clause); the paper's sentence ``the projection step itself is not [machine-checked]'' is accurate.

\subsection{Theorem 2: Hermite attenuation}

$\E_{P_t}[H_\nu(x_t)]=e^{-|\nu|t}\E_{P_0}[H_\nu(a)]$ for $\E|a|^{|\nu|}<\infty$. The generating-function proof is standard and complete; coordinates of $z$ are independent given $a$, so the product form holds pointwise in $a$; Fubini needs $\E|H_\nu(a)|<\infty$, which finite $|\nu|$-th moments give. \sev{verified}. Two remarks the paper should absorb:
\begin{itemize}
\item The cumulant consequence $\kappa_m(x_t)=e^{-mt}\kappa_m(a)$, $m\ge3$, follows directly from $x_t=\alpha a+\text{Gaussian}$ and needs no Hermite polynomials; the Hermite statement is the natural one for \emph{moments}, which is what a Fisher information is built from.
\item The \hi{dictionary} that the abstract currently blurs. For a feature whose lowest Hermite footprint has degree $m$: amplitude of the moment, $e^{-mt}$ (exact, all $t$, Theorem~2); Fisher information, $e^{-2mt}$ (asymptotic, Theorem~3/\ref{thm:general}); sample complexity for fixed local precision, $e^{2mt}$ (asymptotic in $M$, Corollary~\ref{cor:lam}). Correlation: $m=2$, so $e^{-2t}$, $e^{-4t}$, $e^{4t}$. Fourth cumulant: $m=4$, so $e^{-4t}$, $e^{-8t}$, $e^{8t}$. The squared score gap is the square of a degree-4 amplitude: $e^{-8t}$. The abstract writes ``correlation survives as $e^{-4t}$'' (information) two lines after ``attenuates every Hermite mode of degree $m$ by $e^{-mt}$'' (amplitude); one sentence stating the dictionary removes the ambiguity everywhere.
\end{itemize}

\subsection{Theorem 3: score gap and Fisher information}
\label{sec:thm3}

\paragraph{Exact identities.} $P_t(x)=\gamma_\Delta(x)\E_{P_0}\exp(s\cdot a-b|a|^2/2)$ with $s=\alpha x/\Delta$, $b=\alpha^2/\Delta$; $\E[a\mid x]=\nabla_s\Lambda$; $S=-x/\Delta+(\alpha/\Delta)\nabla_s\Lambda$. All correct and require only a density for $P_0$.

\paragraph{The gap in the proof.} Two statements in Appendix A.3 are false under the stated hypothesis.
\begin{enumerate}
\item \emph{``$\Lambda$ is real-analytic in $(s,b)$ on $\R^L\times(0,\infty)$, and its Taylor coefficients at the origin are the joint cumulants''.} $\Lambda$ is analytic in $s$ for every fixed $b>0$ (the Gaussian damping makes the tilted measure integrable), but the expansion is taken at $b=0$, and $\Lambda(s,0)=\log\E e^{s\cdot a}$ is $+\infty$ for every $s\ne0$ when $P_0$ has only finitely many moments (Student-$t$; any regularly varying tail). The Taylor coefficients exist up to the order of the finite moments, the function is $C^k$ but not analytic, and a Lagrange remainder needs the derivatives at intermediate $(\theta s,\theta b)$, i.e.\ moments of the tilted measure $\propto P_0(a)\exp(-\theta\alpha^2|a-x/\alpha|^2/2\Delta)$, which on the bulk $\{|\alpha x|\le1\}$ are bounded only by $C_k e^{\theta|x|^2/2\Delta}$---not uniformly integrable against $P_t$ at the order claimed.
\item \emph{``$P_t(|\alpha x|>1)\le\exp(-c/\alpha^2)$ since $x_t$ is a Gaussian convolution''.} The event $\{|\alpha x|>1\}$ contains $\{\alpha^2|a|>1/2\}$, whose probability is $\Theta(\alpha^{2\nu})$ for a tail index $\nu$ and $\le\E|a|^8(2\alpha^2)^8=O(\alpha^{16})$ under eighth moments---polynomial, not exponential. The conclusion survives (an $O(\alpha^{16})$ tail times a polynomial crude bound is $o(\alpha^8)$), but the sentence is wrong and a referee will not read past it.
\end{enumerate}

\paragraph{Is the theorem true under finite eighth moments?} Probably yes for the score gap, and the eighth moment is close to sharp, but the present proof does not show it and this audit could not close it cleanly. Heuristic for the scalar case with $P_0(a)\sim c|a|^{-\nu-1}$: on $\{|x|>K\}$ the true posterior mean is $\approx x/\alpha$ (the likelihood dominates a polynomial prior tail), so $S(x)\approx\partial\log P_0(x/\alpha)/\alpha\approx-(\nu+1)/x$ while $\Sg(x)\approx-x$, and the tail contributes $\int_{|x|>K}x^2P_t(x)\,dx\asymp\alpha^\nu K^{2-\nu}$ to the squared gap. This is $o(\alpha^8)$ iff $\nu>8$: finite eighth moments are exactly the borderline, and at $\nu\le8$ the tail term $e^{-\nu t}$ is leading. So Theorem~3's hypothesis is the right \emph{condition} for the conclusion but not a hypothesis under which the given \emph{proof} works.

\paragraph{The repair.} State the theorem under
\begin{assumption}[E]
$\E_{P_0}e^{\lambda|a|}<\infty$ for some $\lambda>0$.
\end{assumption}
Then: (i) $\Lambda(s,b)$ is finite and real-analytic in $s$ on $|s|<\lambda$ jointly with $b\ge0$ (the integrand is dominated by $e^{\lambda|a|}$), so the Taylor expansion at $(0,0)$ with a Lagrange remainder is legitimate on $\{|s|\le\lambda/2\}$, i.e.\ on the bulk $\{|x|\le\lambda\Delta/2\alpha\}$, with all tilted moments bounded uniformly there; (ii) the tail $\{|x|>\lambda\Delta/2\alpha\}\subset\{\alpha|a|>c/\alpha\}\cup\{\sqrt\Delta|z|>c/\alpha\}$ has probability $\le e^{-\lambda c/\alpha^2}\E e^{\lambda|a|}+e^{-c^2/2\alpha^2}$, which \emph{is} $\exp(-c'/\alpha^2)$; (iii) on the tail $|S-\Sg|^2$ is bounded by $C(1+|x|^2+\E[|a|^2\mid x])/\Delta^2$ and $\E[\E[|a|^2\mid x]\mathbf 1_{\text{tail}}]=\E[|a|^2\mathbf 1_{\text{tail}}]$ by the tower property, so Cauchy--Schwarz with the exponential tail probability makes the tail contribution below every power of $\alpha$. The rest of Appendix A.3 (grading, Wick, the leading terms) is correct as written. Every family in the paper's laboratory except Student-$t$ satisfies (E); Student-$t$ is the family whose behaviour the theorem does not cover and Section~6.2 shows it does not follow the theorem beyond $t\approx4$. The sentence ``Student-$t_5$, whose sixth moment is infinite, departs as the moment condition predicts'' becomes ``Student-$t_5$ violates Assumption~(E); Proposition~\ref{prop:heavy} gives its rate''.

\paragraph{Constants.} Tensor conventions: $\Frob{\kappa_4}^2=\sum_{ijkl}\kappa_4[ijkl]^2$ over ordered index tuples; the code (\texttt{theory.score\_gap\_constant}) sums the full $(L,L,L,L)$ tensor, and the chain measurement reproduces $\tfrac16\Frob{\kappa_4}^2$ to $1$--$6\%$, so the convention is consistent. Wick multiplicities, for symmetric $T$:
\[
\E_\gamma\Bigl[\bigl(\tfrac16\textstyle\sum_{jkl}T_{mjkl}{:}x_jx_kx_l{:}\bigr)^2\Bigr]=\tfrac{3!}{36}\sum_{jkl}T_{mjkl}^2,
\]
\[
\E_\gamma\Bigl[\bigl(\tfrac1{24}\textstyle\sum T_{ijkl}{:}x_ix_jx_kx_l{:}\bigr)^2\Bigr]=\tfrac{4!}{24^2}\Frob{T}^2=\tfrac1{24}\Frob{T}^2,
\]
giving $\tfrac16\Frob{\kappa_4}^2$ for the score gap and $\tfrac1{24}\Frob{T}^2$ for the information. Both checked by hand and, for the shape constant, by scalar quadrature to $0.25\%$ (Section~6.5). Hermite normalisation: probabilists', $\E_\gamma\Hen{n}^2=n!$; consistent throughout.

\paragraph{Nuisance projection.} The efficient information of a degree-4 coordinate after eliminating degree-2 nuisances is $I_{44}-I_{42}I_{22}^{-1}I_{24}$ with $I_{42}=O(\alpha^8)$ (the $\alpha^6$ term vanishes by $\gamma$-orthogonality), $I_{22}=\Theta(\alpha^4)$, so the correction is $O(\alpha^{12})\ll\alpha^8$. Correct. But the manuscript should say \emph{which} nuisances: the argument covers nuisances of lower Hermite degree. A nuisance of the \emph{same} degree (another shape coordinate, as in the $C=8$ mixture) reduces the efficient information by an $O(1)$ factor, which is one reason the one-parameter bound of Section~4 is only a reference.

\paragraph{Remainder and crossover.} $\epsilon_t=O(e^{-2t}\lmax)$ is proved only for the Gaussian chain (where $\Sigma_t^{-1}=I-\alpha^2(\Sigma_0-I)+\dots$ makes the correction explicit). For non-Gaussian $P_0$ the proof gives $\epsilon_t=O(e^{-2t})$ with a constant depending on $P_0$ through its moments and $L$; the identification of that constant with $\lmax$ is a heuristic. $\tstar=\tfrac12\log\lmax$ therefore has theorem-level meaning for the Gaussian chain and is the natural scale otherwise. There is a second, exact meaning worth stating: for the Gaussian chain $P_t/\gamma\in L^2(\gamma)$ iff $\lmax(\Sigma_t)<2$ iff $e^{2t}>\lmax(\Sigma_0)-1$, so $\tfrac12\log(\lmax-1)$ is exactly the time after which the Hermite expansion of $P_t/\gamma$ converges in $L^2(\gamma)$; for a Laplace $P_0$ it never does, which is why the proof must go through $\Lambda$.

\paragraph{Uniformity.} The statement is for fixed $P_0$ (fixed $L$) as $t\to\infty$; nothing is claimed uniformly in $L$, and the constants grow with $L$ ($\Frob{\kappa_4}^2\propto L$). The manuscript says ``as $t\to\infty$''; it should add ``for fixed $L$''.

\paragraph{Verdict.} \sev{supported but overstated}: true and verified under Assumption~(E); the stated hypothesis is not the one the proof uses; the tail sentence is false as written; the crossover is exact only for the Gaussian chain.

\subsection{General order: from a remark to a theorem}

\begin{theorem}[Order-by-order information]
\label{thm:general}
Let $P_0^\theta$ be a family satisfying Assumption~(E) uniformly near $\theta_0$, differentiable in $\theta$ in the sense that $\partial_\theta\E_{P_0^\theta}[H_\nu(a)]$ exists for every $\nu$. Let $m$ be the smallest degree with $\partial_\theta\E[H_\nu(a)]\ne0$ for some $|\nu|=m$. Then, for fixed $L$, as $t\to\infty$,
\[
  I_t(\theta)=e^{-2mt}\sum_{|\nu|=m}\frac{\bigl(\partial_\theta\E_{P_0}[H_\nu(a)]\bigr)^2}{\nu!}\,\bigl(1+O(e^{-2t})\bigr).
\]
\end{theorem}
\noindent\emph{Proof sketch.} Under (E) write $P_t/\gamma_\Delta=e^{\Lambda(s,b)}$ and grade as in Appendix A.3. The two formal series in $\alpha$ for $\log(P_t/\gamma)$---the graded expansion of $\Lambda$ and $\log\bigl(1+\sum_{|\nu|\ge2}\alpha^{|\nu|}\E[H_\nu]H_\nu(x)/\nu!\bigr)$, the latter from Mehler's kernel---agree coefficient by coefficient, because both are polynomial identities in the moments of $P_0$ up to each order (they can be compared on any compactly supported $P_0$, where both converge, and a polynomial identity in finitely many moments that holds on that class holds identically). The grade-$m$ coefficient of $\partial_\theta\log(P_t/\gamma)$ is $\alpha^m\sum_{|\nu|=m}\partial_\theta\E[H_\nu]H_\nu(x)/\nu!$ plus $\partial_\theta$ of products of lower-grade terms, each of which contains a factor with vanishing $\theta$-derivative. Its $\gamma$-second moment is $\alpha^{2m}\sum_{|\nu|=m}(\partial_\theta\E H_\nu)^2/\nu!$ by orthogonality; replacing $\gamma$ by $P_t$ and adding the remainder costs $1+O(\alpha^2)$ exactly as in the $m=4$ case, with the tail handled by (E). $\square$

\noindent\emph{Consistency.} $m=2$: $\sum_i(\partial\Sigma_{ii})^2/2+\sum_{i<j}(\partial\Sigma_{ij})^2=\tfrac12\Frob{\partial\Sigma_0}^2$, eq.~(4). $m=4$ with $\partial\Sigma_0=0$: the degree-4 Hermite moments differ from the twin's by $\kappa_4$ (Theorem~2's consequence), the multiset $\nu$ collects $4!/\nu!$ ordered tuples, so $\sum_\nu(\partial\kappa_4[\nu])^2/\nu!=\Frob{\partial\kappa_4}^2/24$, eq.~(5). The constant is the \emph{Hermite norm} of the derivative of the degree-$m$ moment tensor; for $m=6$ it is $\Frob{\partial\kappa_6+\text{(products)}}^2/6!$, where the products of lower cumulants do not vanish unless all lower-degree derivatives vanish---which is the hypothesis. Identifiable tangent directions are exactly those with some $m<\infty$; a direction with $\partial_\theta\E H_\nu=0$ for all $\nu$ is not identifiable at any noise level, which for a law determined by its moments (implied by (E)) means $\partial_\theta P_0=0$.

\subsection{Sample complexity: what a Fisher information licenses}

\begin{corollary}[Local asymptotic minimax, order by order]
\label{cor:lam}
Fix $t$ and let $I_t^{\mathrm{eff}}(\theta)$ be the efficient Fisher information of one sequence about a coordinate $\theta$ whose lowest footprint has degree $m$. For $M$ i.i.d.\ sequences observed at level $t$ and any estimator sequence $\hat\theta_M$,
\[
 \liminf_{M\to\infty}\ \sup_{|\theta'-\theta|\le c/\sqrt M}\ M\,\E_{\theta'}(\hat\theta_M-\theta')^2\ \ge\ \frac{1}{I_t^{\mathrm{eff}}(\theta)}
 \;=\;\frac{e^{2mt}}{C_m(\theta)}\bigl(1+O(e^{-2t})\bigr),
\]
and the maximum-likelihood estimator attains the bound. In words: for fixed local precision the number of sequences \emph{necessary and sufficient} grows as $e^{2mt}$, with the constant of Theorem~\ref{thm:general}.
\end{corollary}
\noindent This is H\'ajek--Le Cam for a regular parametric model ($P_t^\theta$ is smooth in $\theta$ with finite, non-degenerate information at every $t>0$; differentiability in quadratic mean follows from the smoothness of $\Lambda$ in $\theta$ under (E)), and the MLE's efficiency is the standard consequence. What it does \emph{not} say, and the manuscript must not: (i) anything uniform in $t$ (the LAN regime needs $M\gg1/I_t$, i.e.\ $M\gg e^{2mt}$, so the two limits do not commute); (ii) anything about absolute precision for a bounded coordinate---once $1/\sqrt{MI_t}$ exceeds the coordinate's range the local statement is empty, which is exactly the collapse E06 observes at $t\approx1$; (iii) anything about EM at a finite iteration budget. The ratio statement ``$e^{4t}$ fewer effective samples for shape than for covariance'' is the ratio of two such bounds and inherits the same regime. ``Requires'' in Corollary~5 should read ``needs, in the local asymptotic minimax sense''.

\subsection{Heavy tails: outside the theorem, with a rate of their own}

\begin{proposition}[Heavy-tailed regime; supported numerically, proof sketched]
\label{prop:heavy}
Let $P_0$ be a chain whose innovation density is regularly varying with index $\nu>4$, $\varphi(\varepsilon)\sim c|\varepsilon|^{-\nu-1}$. Then $\E_{P_t}\|S-\Sg\|^2=\Theta(e^{-\min(\nu,8)t})$: the Hermite term $\tfrac16\Frob{\kappa_4}^2e^{-8t}$ leads while $\tfrac16\Frob{\kappa_4}^2e^{-8t}\gg c'Le^{-\nu t}$ and the tail term leads afterwards. For $\nu>8$ the theorem's conclusion holds; for $\nu\le8$ the moment condition of Theorem~3 is violated and the asymptotic rate is $e^{-\nu t}$.
\end{proposition}
\noindent\emph{Evidence.} Scalar quadrature on an untruncated domain (Section~6.2, running as this is written) and the chain at $(N_g,A)=(1601,16)$: for $\nu=5$ the ratio of the measured gap to $\tfrac16\Frob{\kappa_4}^2e^{-8t}$ is $0.57$, $0.99$, $18$ at $t=3,4,5$ and the local slope moves from $-7.7$ to $-4.5$ between $t=3.5$ and $t=5$. The crossover is late because $\Frob{\kappa_4}^2$ is large ($452$ per chain for $\nu=5$) and the tail constant is small. The $e^{-8t}$ constant is therefore \emph{reached} by Student-$t_5$, which the manuscript currently denies, and then abandoned, which it currently misattributes. The fourth-order Edgeworth correction of Section~6.6 closes the gap on every (E)-family and fails on Student-$t$ at exactly the levels where the tail term leads ($82\%$ closed at $t=2.1$, $8\%$ at $t=3.5$), a second, independent signature of the same mechanism.

\subsection{Locality: the exact radius-$r$ floor and its Gaussian closed form}

\begin{proposition}[Influence radius of the Gaussian score]
\label{prop:radius}
For the infinite Gaussian AR(1) chain, $(\Sigma_t^{-1})_{i,i+k}\propto r(\rho,t)^{|k|}$ for $|k|\ge2$, with
\[
 r(\rho,t)=\frac{A-\sqrt{A^2-B^2}}{B},\qquad A=\Delta(1+\rho^2)+\alpha^2(1-\rho^2),\quad B=2\Delta\rho .
\]
$r\to0$ as $t\to0$ (the score is local) and $r\to\rho$ as $t\to\infty$ (the score's off-diagonal part is $-\alpha^2(\Sigma_0-I)$, with the prior's correlation length). The best radius-$r$ linear estimator's excess error over the LMMSE is the energy of the coefficients beyond $r$, computable in closed form; the best radius-$r$ estimator of any kind is $\E[a_i\mid x_{i-r},\dots,x_{i+r}]$, computable exactly by sum-product on the sub-chain with the exact marginal of its first site as initial factor (Section~6.4).
\end{proposition}
\noindent\emph{Proof.} The symbol of $\Sigma_t$ is $g(\omega)=\alpha^2(1-\rho^2)/(1-2\rho\cos\omega+\rho^2)+\Delta$; $1/g=(1-2\rho\cos\omega+\rho^2)/(A-B\cos\omega)$, and the Fourier coefficients of $1/(A-B\cos\omega)$ are $r^{|k|}/\sqrt{A^2-B^2}$ with $r$ the root of $r+1/r=2A/B$ inside the unit disc; the numerator is a trigonometric polynomial of degree one and only modifies $|k|\le1$. $\square$

\noindent This is the analytical form of ``the window head plateaus''---\emph{if} the plateau is where the floor is. Section~6.4 measures the floor; the verdict is in Section~4.

\subsection{Length extrapolation}

The EM--BP estimator is length-equivariant by construction: the fitted kernel is a function on $\R^2$, and the recursion at length $L'$ is the same recursion. Its relative error at length $L'$ is a functional of the kernel error alone; since the chain is stationary away from the first site and the posterior has finite correlation length (Proposition~\ref{prop:radius} for the Gaussian part), the per-site error converges as $L'\to\infty$, which is why $0.023$ is flat from $32$ to $256$. A convolutional or recurrent network is equivariant up to its receptive field or state; a transformer with a clipped relative-position bias ($\max\_dist=32$) is equivariant only up to that clip, and at $L'>64$ sees position pairs it never saw. This is a remark, not a theorem, and the manuscript treats it correctly as one.

\subsection{Misspecification: a perturbation theory of where structure stops helping}

Let the truth be Gaussian with covariance $\Sigma(\varepsilon)$, $\Sigma(0)=\Sigma_0$ an AR(1) covariance, and let the fitted family be Gaussian AR(1) (the $C=8$ mixture fits a Gaussian innovation to a Gaussian truth: E08's fitted excess kurtosis is $0.04$--$0.34$). The population EM limit under the one-view protocol is the KL projection
\[
 \theta^\ast(\varepsilon)=\arg\min_{(\rho,q)}\sum_{t\in\mathcal T}\mathrm{KL}\bigl(N(0,\Sigma_t(\varepsilon))\,\|\,N(0,\hat\Sigma_t(\rho,q))\bigr),
\]
and the asymptotic ($M\to\infty$) relative score error is
\[
 b(\varepsilon)^2=\frac{\sum_t\operatorname{tr}\bigl[(\hat\Sigma_t^{-1}-\Sigma_t^{-1})\Sigma_t(\hat\Sigma_t^{-1}-\Sigma_t^{-1})\bigr]}{\sum_t\operatorname{tr}\Sigma_t^{-1}}\Bigg|_{\theta^\ast(\varepsilon)} .
\]
\begin{proposition}[Order of the bias]
\label{prop:bias}
$b(\varepsilon)=\Theta(\varepsilon)$ whenever $\partial_\varepsilon\Sigma(0)$ is not in the tangent space of the AR(1) family, and $b(\varepsilon)=\Theta(\varepsilon^2)$ when the perturbation enters the law only through $\varepsilon^2$ (a perturbation odd in a latent variable). The finite-$M$ risk is $\sqrt{b(\varepsilon)^2+v/M}$ to leading order, with $v$ the well-specified variance constant.
\end{proposition}
\noindent Section~6.3 computes $b(\varepsilon)$ for the four E08 mechanisms with no Monte Carlo: it reproduces the $M=2048$ EM--BP error to $\pm0.003$ at every strength of every mechanism, gives $b\sim0.60\,\varepsilon$ (AR(2)), $1.75\,\varepsilon$ (site-varying $\rho$), $1.07\,\varepsilon^{0.8\to1}$ (long-range), and $0.15\,\beta^2$ (rank-one latent). The rank-one latent is the only even perturbation ($\Sigma=(\Sigma_y+\beta^2\mathbf 1\mathbf 1^\top)/(1+\beta^2)$), and it perturbs one eigendirection of $L=32$, so its effect on the score energy is $O(1/L)$ as well as $O(\beta^2)$: two reasons, both quantitative, why it is ``absorbed''. This replaces the manuscript's qualitative ``because it does not break the conditional independence'' (which is false: the latent does break Markovianity) with the right reason.

%% ------------------------------------------------------------------------------------
\section{Experimental audit}
\label{sec:exp}

For every experiment: the question, the varied factor, what is frozen, pairing, the inferential unit, selection, uncertainty, the estimand, and the four confounds (budget, numerics, model selection, generalisation beyond the point). Verdicts use the brief's vocabulary.

\subsection{E01 --- order-by-order decay (theory verification)}
\hi{Question.} Do the exact score and the exact information curves follow Theorem~3's rates and constants? \hi{Varied.} $t$ (40 fine levels for the Gaussian chain; twelve plus $3.5, 4$ otherwise), $\rho$, $\beta$, family. \hi{Frozen.} $L=32$, grid $(801,8)$ for the non-Gaussian parts. \hi{Unit.} Replicates (8 for the shape information, 4 for the closure gap) with common noise across $t$. \hi{Selection.} None; slopes are fitted by a declared tool on declared windows.

\emph{Gaussian part.} Exact, no Monte Carlo: slope $-4.00$ and scaled information $1.000$ at every $\rho$; the approach to the asymptote is governed by $\tstar$. \sev{verified}.

\emph{Shape information.} Slopes $-7.94$ ($\beta=2$), $-8.01$ ($\beta=3$) on the last intervals; the Laplace member floors at $t\approx2.5$. \sev{verified} for the \emph{rate}. The \emph{constant} of eq.~(5) is \emph{not} reproduced on the chain: the plateau of $I_\beta e^{8t}$ is $0.81$ of the prediction at $\beta=2$ and $0.92$ at $\beta=3$ (Section~6.5), while the scalar quadrature of E21 gives $1.000$. The manuscript's ``verifies the rates and the constants'' must name which constants are verified where: the score-gap constant on the chain to $1$--$6\%$, the covariance-information constant on the chain exactly, the shape-information constant in the scalar case to $0.1\%$ and on the chain to $10$--$20\%$ before the outer-product estimator's floor. \sev{supported but overstated}.

\emph{Closure gap.} Laplace and bimodal reach the constant ($0.99$--$1.06$, $0.88$--$0.98$ at $t\ge3$); uniform hits the double-precision floor at $t\approx3.5$; Student-$t_5$ ``departs''. E18 (Section~6.2) shows the departure on the $(801,8)$ grid is truncation of the prior: at $(3201,32)$ the same $512$ Student-$t_5$ sequences follow slopes $-7.8$ to $-7.9$ through $t=5$, and with $8192$ sequences the constant is $1.03$ $[0.99,1.09]$ at $t=4$ (Section~6.3); $\nu=7, 9, 12$ likewise. \sev{confounded} (numerics) for the Student-$t$ curve; the panel must be redrawn from E18.

\subsection{E02a --- capacity}
\hi{Question.} Does mixture capacity buy score accuracy beyond the covariance? \hi{Varied.} $C\in\{1,2,4,8,16\}$, $M$. \hi{Frozen.} Bundles (hashed), initialisation stream, ladder to $1200$, selection rule. \hi{Paired.} Yes, within (seed, $M$) on identical test bundles. \hi{Unit.} Sixteen seeds. \hi{Selection.} Checkpoint on a validation bundle scored against the true kernel (an oracle, shared by every arm). \hi{Budget.} $C\ge4$ at $M=2048$ selects the last rung in $56$--$69\%$ of cells, but E02a\_v1 (ladder to $400$) reproduces every overlapping cell to $0.001$: the estimates have plateaued (the validation curve moves $0.0147\to0.0145$ over the last $400$ iterations). \sev{verified}. Two things to add: the best capacity rises with $M$ ($C=2$ at $M=32$, $C=16$ at $M=2048$), which the text says; and the ``$M=\infty$'' column is a KL projection, not a score-optimal fit, which the text also says. The fitted excess kurtosis at $M=32$ ranges $0.3$--$8.9$ across seeds---the shape is barely identifiable at $M=32$, and $C=8$'s advantage there ($0.099$ vs $0.174$) is mostly covariance-plus-a-little-shape; worth one clause.

\subsection{E02b --- the family}
Six members, $M=512$, $C\in\{1,8\}$, sixteen seeds. The $C=1$ arm traces the LMMSE floor's $V$ within its s.e.\ at every member; $C=8$ stays at $0.017$--$0.035$; the gain reverses by $-0.008$ $[-0.011,-0.005]$ at the Gaussian member. \sev{verified}, and the cleanest single figure in the paper. One caveat the caption should carry: at the Gaussian member the floor is exactly zero, so the reversal is ``extra components cost estimation variance'', measured against an exact score, not against a floor.

\subsection{E03 --- baselines, matched interface, length}
\hi{Screen.} Two development seeds, one run of 16k steps per configuration, 68 configurations, ranked by pooled bulk risk. Adequate for ranking; the confirmatory protocol is separate. The transformer wins its family at lr $0.003$ while every other family wins at $0.001$---fine, the screen is per family.

\hi{Confirm.} Sixteen seeds, four sizes, 32k steps, both parameterisations screened, checkpoint selected on the oracle validation bundle. At $M=2048$ the MLP and the transformer select the last rung in $88\%$ of cells and the bi-GRU in $31\%$: the networks with the most headroom were still improving at 32k steps. \emph{The $2.83\times$ paired factor against the best network is therefore an upper bound on the true advantage under the standard recipe}, exactly as the manuscript says of EM--BP's own at-cap rate in the capacity experiment and does not say here. \sev{supported but overstated} until that sentence is added; the direction is not in doubt (a factor of three at sixteen seeds, and the bi-GRU that carries the comparison is at cap in only a third of cells).

\hi{Matched.} Two arms in \texttt{matched.csv}: \texttt{ism\_one\_view} (implicit score matching on $x_t$ only; checkpoint selected on the ISM objective of a held-out noisy bundle; no clean chain, no true score anywhere) and \texttt{dsm\_one\_corruption} (denoising score matching on $(a,x_t)$ pairs with one fixed corruption; checkpoint selected on the oracle bundle). The generator (\texttt{make\_baseline\_numbers.py}) takes the \emph{best} of the two by risk, which at every size is the DSM arm, and labels it ``the comparison at matched information''. It is not: the DSM arm sees the clean sequence. Table~\ref{tab:matched} gives both, paired with EM--BP $C=8$ on identical bundles.

\begin{table}[h]\centering\small
\caption{The two ``matched'' arms against EM--BP ($C=8$), all-site pooled error, sixteen seeds, identical test bundles; paired factor = mean over seeds of (network/EM--BP), percentile bootstrap. Source: \texttt{outputs/E03/matched.csv}, \texttt{outputs/E02a/capacity.csv}.}
\label{tab:matched}
\begin{tabular}{@{}llrrrrl@{}}
\toprule
arm & what the network sees & $M$ & network & EM--BP & paired factor & wins \\
\midrule
ISM, one view & $x_t$ only (matched) & 2048 & 0.165 & 0.023 & 7.9 [6.6, 9.3] & 16/16 \\
 & & 512 & 0.211 & 0.036 & 6.0 [5.5, 6.5] & 16/16 \\
 & & 128 & 0.401 & 0.055 & 8.1 [6.8, 9.7] & 16/16 \\
 & & 32 & 1.205 & 0.099 & 13.3 [11.3, 15.7] & 16/16 \\
DSM, one corruption & $(a, x_t)$: the clean chain too & 2048 & 0.155 & 0.023 & 7.3 [6.2, 8.7] & 16/16 \\
 & & 512 & 0.206 & 0.036 & 5.9 [5.4, 6.4] & 16/16 \\
 & & 128 & 0.272 & 0.055 & 5.5 [4.6, 6.7] & 16/16 \\
 & & 32 & 0.475 & 0.099 & 5.3 [4.5, 6.2] & 16/16 \\
\bottomrule
\end{tabular}
\end{table}

Classification against the brief: the standard recipe is interface~A; the fixed-corruption arm is interface~B (clean target, no fresh corruption)---strictly more information than EM--BP, less augmentation than A; the ISM arm is interface~C, genuinely observation-matched. The headline must be C. It happens to be larger, and it is the only arm in the paper selected without an oracle. The ISM window head at $M=32$ diverges (risk $1.2$, the ISM objective on the held-out bundle rising from $20$ to $2500$ across the ladder)---known behaviour of implicit score matching on tiny samples---so the $M=32$ ISM number is not a measurement of the architecture; report $M\ge128$ or say why. \sev{contradicted} as worded; \sev{verified} once reworded around C.

\hi{Length.} EM--BP $0.023$ at every length; window head $0.118\to0.082$; bi-GRU/TCN/U-Net improve slightly; transformer $0.063\to0.703$. \sev{verified}. Note that the window head improves \emph{more} than the global architectures with length, which says a sizeable part of its $L=32$ error is at the boundary (zero padding); the ``plateau'' should be read on the bulk region as well (Section~6.4).

\subsection{E05 --- generation}
Sixteen seeds, common Brownian path within a seed, 400 steps with the discretisation measured separately ($0.019\to0.0033$ strong error from 100 to 400 steps). Deficit of the fitted kernel $+0.09$ nats $[+0.00,+0.19]$; kurtosis $1.80$ vs $1.94$; the covariance twin $+1.70$ with kurtosis $0.00$; the isotropic $+46$. \sev{verified}; the phrase ``small but real'' should be ``$0.09$ nats, interval touching zero''. The report and the paper now agree (the audit tool caught the stale four-seed table).

\subsection{E06 --- learning through noise}
\hi{Question.} How do the errors of $\hat\rho$ and $\hat\kappa_4$ grow with the observation level? \hi{Varied.} $t$. \hi{Frozen.} $M=512$, $C=8$, ladder to $1200$, one view per sequence. \hi{Unit.} Sixteen seeds per level. \hi{Selection.} Oracle validation risk. \hi{Estimand.} RMSE over seeds of the fitted coordinate---declared. What the data say per level (Table~\ref{tab:e06}):

\begin{table}[h]\centering\small
\caption{E06 per level: collapse fraction (fitted excess kurtosis $<1$), dispersion and bias of the non-collapsed fits, the one-parameter Cram\'er--Rao reference read from E01, and the fraction of cells whose selected checkpoint is the last rung. Source: \texttt{outputs/E06/levels.csv}; reference from \texttt{outputs/E01/shape\_information.csv}.}
\label{tab:e06}
\begin{tabular}{@{}rrrrrrrrr@{}}
\toprule
$t$ & collapsed & sd (all) & sd (good) & bias (all) & CRB & RMSE/CRB & sd(good)/CRB & at cap \\
\midrule
0.05 & 0/16 & 0.29 & 0.29 & $-0.03$ & 0.24 & 1.16 & 1.19 & 0.25 \\
0.10 & 0/16 & 0.48 & 0.48 & $+0.14$ & 0.32 & 1.51 & 1.49 & 0.25 \\
0.20 & 0/16 & 0.34 & 0.34 & $-0.07$ & 0.47 & 0.71 & 0.72 & 0.38 \\
0.40 & 0/16 & 0.69 & 0.69 & $-0.29$ & 0.84 & 0.87 & 0.83 & 0.25 \\
0.70 & 4/16 & 1.46 & 0.76 & $-0.66$ & 1.79 & 0.87 & 0.42 & 0.31 \\
1.00 & 15/16 & 0.98 & --- & $-2.63$ & 3.81 & 0.73 & --- & 0.12 \\
1.40 & 16/16 & 0.10 & --- & $-2.92$ & 12.6 & 0.23 & --- & 0.88 \\
2.00 & 16/16 & 0.13 & --- & $-3.05$ & 105 & 0.03 & --- & 0.75 \\
\bottomrule
\end{tabular}
\end{table}

The ``$1.02\times$'' is the mean of the RMSE/CRB column over the rows with RMSE $<2$: $1.16$, $1.51$, $0.71$, $0.87$, $0.87$. Three of those five are \emph{below} one, which for an estimator compared with an unbiased-estimator bound means the estimator is biased (it is: $-0.29$ and $-0.66$) or the reference is for the wrong family (it is: one parameter against a mixture). At $t=0.7$ the surviving cluster has dispersion $0.42$ of the bound and the collapsed quarter sits at the range. The estimator is not ``at the bound''; it is a shrinkage estimator that shrinks toward a Gaussian innovation as the information falls, and the validation-selected early stopping is the shrinkage. \sev{contradicted} as worded. What is \sev{verified}: for $t\le0.4$ the fits are unimodal and their dispersion is within $\pm20\%$ (the sampling error of a 16-seed sd) of the one-parameter reference; collapse starts at $t=0.7$ and is total at $t=1$, one level below the level at which the reference bound crosses the coordinate's range ($0.98$); $\hat\rho$ is within $0.01$ through $t=1$ and its dispersion is at the reference at every level up to $1.4$.

\hi{Budget.} At $t\ge1.4$, $75$--$88\%$ of cells select the last rung; flagged correctly; E06c is running (eleven cells in). Its first cells already answer the question in the direction the flag anticipated: at $t=2$, seed 0's fit moves from $\hat\rho=0.577$ at rung $220$ of the 1200-ladder to $\hat\rho=0.800$ at rung $3746$ of the 8000-ladder (truth $0.85$); at $t=1.4$, seed 1 moves from $0.831$ to $0.844$; the collapsed shape stays collapsed. The ``leaves the bound'' clause was budget, at least in part, and the manuscript's decision to report it as unresolved was right. \sev{pending} (frozen numbers when the run lands).

\hi{Slopes.} No asymptotic window exists at $\rho=0.85$ with $M=512$: the reference bound crosses the range at $t=0.98$ and $\tstar=1.17$. Pre-asymptotic slopes $1.29$ and $3.03$ on $[0.2,0.7]$ are reported as such. \sev{verified}; the introduction's ``confirms'' is not.

\subsection{E07 --- the family ladder; E11 --- the envelope}
E07 is deterministic (two replicates of 512 chains). Its Student-$t$ row at $\rho=0$ is flagged because E11 reports ``belief mass in the outer $2\%$ = 1.00''. Section~6.1 reproduces E11's chains: seeds 0--3 have one, one, one and two sites beyond $|a|=8$ among $8192$ each (analytic expectation $1.2$), seed 2's at $|a|=15.4$---beyond the $A=12$ reference too, which is why that seed's ``drift'' at $(801,12)$ is $2.6\times10^{-5}$ while the others' is $3\times10^{-8}$: both grids truncate the same site. The flat-in-$N_g$ drift is exactly the signature of a handful of pinned sites; the pooled relative error at the lowest level is dominated by them because at $t=0.05$ a site pinned at $8$ instead of $15$ carries a score error of $(\alpha/\Delta)\times7\approx70$. The manuscript's mechanism (truncation, not quadrature) is right; its description (``essentially all of its belief mass'') is a misreading of a max statistic. \sev{supported but overstated}. The right diagnostic is the exceedance count $M L\,P(|a|>A)$, which for Student-$t_5$ at $A=8$ is $1.2$ per $8192$ sites and at $A=16$ is $0.04$; E07b ($A=12$) will still see one site in some replicate, E07c ($A=16$) will usually not. The uniform family's $0.115$ worst-level drift at $\rho=0.5$ is a different mechanism (a discontinuous kernel under trapezoidal quadrature, $O(h)$ not $O(h^2)$) and is not mentioned anywhere; it should be, in one sentence.

\subsection{E08 --- misspecification}
Sixteen seeds, two sizes, four Gaussian mechanisms with exact scores. Verdicts: rank-one latent never reverses (ratio $15.6\to4.0$); AR(2) reverses in $(0.3,0.4)$; site-varying $\rho$ in $(0.10,0.14)$; long-range a tie at $0.4$. \sev{verified}. Two corrections: (i) the ``window head'' is a radius-2, width-64 head trained 16k steps, not E03's screened radius-4 winner---the claim matrix says ``screened'', the text says ``the window head''; name it; (ii) the explanation of the latent's absorption (``does not break the conditional independence'') is wrong---it does break Markovianity---and Section~6.3 gives the right two reasons (even in $\beta$; one direction of 32). With E17 the whole panel gains a theory: EM--BP's error at $M=2048$ is the closed-form bias to $\pm0.003$ at every strength, and the break-even is where $b(\varepsilon)$ meets the network's error.

\subsection{E09, E10 --- capacity references}
E09: the KL-optimal mixture per $C$ on the test bundles ($0.167, 0.040, 0.0058, 0.0022, 0.0007$); the $C\ge8$ fits have a component under two grid cells and are flagged. E10: the same frozen fits in held-out log-evidence per edge; every $C>1$ beats $C=1$ at every $M$, $16/16$. \sev{verified}; together they close the capacity question in both estimands.

\subsection{E12 --- sample complexity}
\sev{confounded} at $t=1$ (at cap $100\%$ at $M\ge4096$, kurtosis non-monotone in $M$), as the project already states. The two one-sided statements stand. E12b is queued with a corrected analysis plan (Section~7).

\subsection{E13 --- regimes}
Fifteen cells, eight seeds. $C=8$ beats $C=1$ everywhere with the interval excluding zero; $C=1$ on the floor to $0.002$ median. \sev{verified}. The $\rho$-trend is reported as a direction, correctly.

\subsection{E14 --- matched comparison across regimes}
Queued (jobs 662797/662798, scheduled $\approx21$ September). Design is sound: ISM for every arm (interface~C throughout, Hutchinson divergence uniformly), networks re-screened per regime, EM--BP on the same views. Eight seeds. The pre-registration predicts ``EM--BP beats the best screened network at every $(\rho,L)$''; nothing in this audit changes that expectation, but the report's decision~7 (wait for it before fixing the framing) is right.

%% ------------------------------------------------------------------------------------
\section{Contradictions between the paper, the report, and the frozen outputs}
\label{sec:contra}

\begin{enumerate}
\item \hi{``Information-matched'' (paper \S5, report \S5.2 and summary item 3).} Both documents call the fixed-corruption DSM arm ``the estimator's own interface''. The frozen file says it trains on $(a,x_t)$ pairs. The report even says ``implicit score matching on that bundle, which uses no clean target at all, lands at $0.165$'' in the same paragraph, i.e.\ it names the matched arm and then quotes the other one as the headline.
\item \hi{``Attains the Cram\'er--Rao bound'' (paper \S4, report \S4 and summary item 2).} The report's ``never worse than $1.51$ for the fourth cumulant'' is true and the ``$1.02$ on average'' is arithmetic on the macro file; neither says that three of the five averaged ratios are below one, that the bound is for a different family, or that the estimator is bimodal from $t=0.7$.
\item \hi{Belief mass outside the domain (paper App.~B, App.~E, Fig.~6; report \S8.2 and summary).} ``Essentially all'' versus one to two sites in $8192$. Same number, wrong reading.
\item \hi{Student-$t_5$ departure (paper Fig.~2c caption and \S3; THEORY\_NOTES \S7).} ``Departs as the moment condition predicts'' versus E18: the constant is reached on a wide domain. The (801, 8) grid, validated at the frozen configuration for the twelve levels, was never validated for the closure panel's $t\le4$ on heavy tails.
\item \hi{The under-resolved fits.} App.~C: ``a generator refuses to emit a table from \ldots cells whose narrowest fitted component is under two grid cells''. Table~1 prints two such cells with daggers and the text says ``1 of 320 fitted cells fail the two-grid-cell resolution gate and are reported, not dropped''. Policy in the code: flagged, not refused. The provenance sentence is the stale one.
\item \hi{Seeds.} \S5: ``Sixteen paired seeds are the inferential unit''. E13 and E14 use eight. The regime paragraph says eight; the protocol paragraph does not qualify.
\item \hi{Which window head.} Claim matrix E18 and E08's registry: ``against the screened window head''; E08's params: radius 2, width 64, 16k steps. E03's screened winner is radius 4, width 64, 32k steps.
\item \hi{\S1 versus \S4.} \S1: ``the experiment confirms it'' ($e^{4t}$); \S4: ``the asymptotic window itself lies beyond the reach of $M=512$''. The latter is right.
\item \hi{Contribution (iii) versus NEGATIVE\_RESULTS.} ``with the predicted $e^{4t}$ separation \ldots measured'' versus ``the predicted exponents are not confirmed by these runs''.
\item \hi{Dangling reference.} \S5g cites \verb|\cref{thm:gap}|, which does not exist; the PDF prints ``??''.
\item \hi{Abstract.} A \texttt{[PENDING]} note in a document labelled ``under review''.
\item \hi{Report summary item 4 versus E03 length.} ``The radius-4 window head \ldots moving only $0.121\to0.118$'' beside a length table in which the same head moves $0.118\to0.082$ as the chain lengthens: the plateau is partly a boundary effect at $L=32$, which neither document says.
\item \hi{Report \S8.1 versus paper \S5g.} The report says ``the mixture beats the Gaussian kernel in all fifteen cells \ldots the margin shrinking smoothly with $\rho$''; the paper says the same and adds ``its trend is not the exponent of \verb|\cref{thm:gap}|''---an honest caveat attached to a broken reference.
\item \hi{Related work.} ``What is not in that literature \ldots is the score-level statement for a Markov (dependent) rather than i.i.d.\ latent structure'': the low-SNR expansions of Prelov--Verd\'u and Guo--Shamai--Verd\'u are for arbitrary input laws, dependent included; the dependence is not the novelty. The constant, the crossover, the exact verification and the learning consequence are (Section~8.3).
\item \hi{Two E06 rows that agree with E06c and one stop reason that does not.} E06c's first four cells reproduce E06's fits exactly (same trajectories); the stop reason changes from ``censored'' to ``converged'' because the longer ladder lets the tolerance fire. Consistent, but the manuscript's ``censored'' counts will change when E06c lands and should be read as ``ladder exhausted'', not ``diverging''.
\end{enumerate}

%% ------------------------------------------------------------------------------------
\section{Additional analyses performed}
\label{sec:analyses}

Everything in this section was run on the laptop during the audit (18--19 September). Each analysis is a registered experiment (E15--E21, \texttt{experiments/registry/}) whose script reproduces it under provenance; numbers are from those runs where complete and from the identical-seed scratch computation otherwise (marked $\dagger$).

\subsection{The Student-$t$ envelope cell is a handful of exceedances (E11 re-read)}
\label{sec:exceed}
E11's chains were regenerated from their tags. Sites with $|a|>8$ among $8192$ per seed: $1, 1, 1, 2$ (seeds 0--3); the analytic expectation for the variance-normalised $t_5$ at $\rho=0$ is $8192\times P(|t_5|>8/\sqrt{0.6})=1.2$. Seed 2's exceedance is at $|a|=15.4$, beyond the $A=12$ reference; seed 3's at $10.0$. Every ``boundary belief mass $=1.00$'' in the frozen file is one of these sites; the drift is flat in $N_g$ because a pinned site is pinned on every grid; and the pooled relative drift at the lowest level ($0.046$ for seed 2) is one site's score error $\approx(\alpha/\Delta)(15.4-8)$ against the energy of $8192$ ordinary sites. At $\rho=0.5$ and $0.85$ no site exceeds $8$ in any seed, yet the ``boundary belief mass'' reaches $0.09$ and $0.17$: sites at $|a|\approx6$--$7$ whose posterior at the lowest level has a tail beyond $8$. Same mechanism, no pinning, drift $\le10^{-2}$.

The correct diagnostic for polynomial tails is the expected exceedance count $ML\,P(|a|>A)$, Poisson-distributed, and the correct statement is that a finite domain cannot make it zero: $A=12$ gives $0.17$ per $8192$ sites, $A=16$ gives $0.04$. A second mechanism in the same table, mentioned nowhere: the uniform family's worst-level drift of $0.115$ at $\rho=0.5$ on $(401,10)$ is the trapezoidal rule on a discontinuous kernel ($O(h)$, not $O(h^2)$), largest at the coarsest spacing.

\subsection{The closure gap on a wide domain (E18, part ``chain'')}
\label{sec:e18}
$\rho=0.85$, 256 chains $\times$ 2, levels $1$--$5$, three grids. Slopes are paired across grids (same chains) and across levels (common noise) and are precise; ratios to the constant from 512 chains are not (next subsection).

\begin{center}\small
\begin{tabular}{@{}llrrrrrr@{}}
\toprule
family & grid & slope $[3,3.5]$ & $[3.5,4]$ & $[4,4.5]$ & $[4.5,5]$ & ratio $t=4$ & ratio $t=5$ \\
\midrule
Student-$t_5$ & $(801,8)$ & $-4.97$ & $-4.21$ & $-4.04$ & $-4.01$ & 27 & 1440 \\
 & $(1601,16)$ & $-7.69$ & $-7.20$ & $-5.73$ & $-4.48$ & 0.99 & 18 \\
 & $(3201,32)$ & $-7.79$ & $-7.90$ & $-7.86$ & $-7.43$ & 0.74 & 1.06 \\
Student-$t_7$ & $(801,8)$ & $-7.71$ & $-6.86$ & $-5.13$ & $-4.21$ & 1.94 & 54 \\
 & $(3201,32)$ & $-7.90$ & $-7.96$ & $-7.98$ & $-7.99$ & 1.08 & 1.09 \\
Student-$t_9$ & $(1601,16)$ & $-7.98$ & $-8.01$ & $-8.02$ & $-8.02$ & 1.04 & 1.02 \\
Laplace & $(801,8)$ & $-7.95$ & $-8.00$ & $-8.02$ & $-8.02$ & 1.59 & 1.56 \\
bimodal & $(3201,32)$ & $-7.87$ & $-7.95$ & $-7.98$ & $-7.99$ & 1.13 & 1.15 \\
uniform & any & $-7.52$ & $-6.25$ & $-4.62$ & $-4.08$ & 3.8 & 148 \\
\bottomrule
\end{tabular}
\end{center}

Three readings. (i) Every departure from $-8$ on the $(801,8)$ grid for the Student-$t$ families is removed by widening the domain: it was the truncated prior in the reference, not the moment condition. (ii) The uniform family's departure is identical on every grid: it is the double-precision floor of the BP posterior means (a squared gap of $\sim10^{-13}$ per chain), which every family reaches near $t\approx5$; the figure shades it. (iii) On a domain wide enough for the sample, Student-$t_5$ follows $e^{-8t}$ through $t=5$ and reaches the constant. The untruncated \emph{scalar} quadrature (part ``scalar'', adaptive domain $3X/\alpha$) tells the other half: for the unit-variance $t_5$ law the local slope of the expected gap is $-5.9, -6.1, -5.9, -5.6, -5.5, -5.4, -5.3, -5.3$ on consecutive half-unit intervals from $t=1$, converging to $-\nu=-5$, and its ratio to $\tfrac16\kappa_4^2e^{-8t}$ grows as $e^{3t}$ ($0.7, 2.0, 5.0, 14, 48, 170, 630, 2400, 9400$ at $t=1,\dots,5$); for $t_7$ the ratio is $1.0$--$1.4$ through $t=5$ (its crossover is later), for $t_9$ it is $1.00$. The expected tail term is carried by sites with $|a|\gtrsim e^{2t}$: at $t=3$ a site contributing to it has probability $\sim10^{-9}$, and $8192\times32$ sites contain none. Proposition~\ref{prop:heavy} is therefore confirmed in expectation and irrelevant in any sample of practical size; the paper says both.

\subsection{The closure constant, and how precisely it is verified (E18, part ``constant'')}
\label{sec:const}
The per-chain squared gap is a degree-six polynomial in nearly Gaussian variables and is heavy-tailed: over $8192$ chains its maximum is $\approx300$ times its median. A 512-chain mean therefore carries a relative standard error of $\approx15\%$: two 512-chain replicates of E01's Laplace closure at $t=3$ give ratios $0.83$ and $1.14$, two 256-chain replicates $1.77$ and $1.33$. E01's ``$0.99$ at $t=3$'' is the mean of four such replicates and its precision is $\approx\pm10\%$, not $1\%$. With $8192$ chains at $(1601,16)$ and a chain bootstrap:

\begin{center}\small
\begin{tabular}{@{}lcc@{}}
\toprule
family & ratio at $t=3$ & ratio at $t=4$ \\
\midrule
Laplace & 0.92 [0.84, 1.01] & 0.99 [0.91, 1.08] \\
bimodal & 0.93 [0.83, 1.05] & 0.98 [0.88, 1.11] \\
uniform & 0.96 [0.89, 1.04] & 3.6 (floor) \\
Student-$t_5$ & 0.70 [0.62, 0.79] & 1.03 [0.99, 1.09] \\
Student-$t_7$ & 0.90 [0.82, 0.98] & 0.96 [0.88, 1.04] \\
Student-$t_9$ & 0.88 [0.81, 0.96] & 0.93 [0.86, 1.01] \\
\bottomrule
\end{tabular}
\end{center}
Every family with an exponential moment is within its interval of $1$ at $t=4$ (the uniform family has hit the double-precision floor there, ratio $3.6$ with a tiny interval---a floor, not a sample); the heavy-tailed members approach from below and $t_5$ is at $1.03$ at $t=4$. The paper quotes these intervals; this does not touch the rates, which are paired and measured to $\pm0.05$.

\subsection{The locality floor (E15)}
\label{sec:e15}
Pooled over the schedule, the error of the exact best radius-$r$ estimator (sum-product on the window with the exact marginal of its first site):

\begin{center}\small
\begin{tabular}{@{}lrrrrrrrrrr@{}}
\toprule
$\rho$ & $r=0$ & $1$ & $2$ & $3$ & $4$ & $6$ & $8$ & $12$ & $16$ & LMMSE (full) \\
\midrule
0.5 & 0.419 & 0.068 & 0.016 & 0.005 & 0.002 & 0.0003 & 0.0000 & 0.0000 & 0.0000 & 0.264 \\
0.85 & 0.666 & 0.216 & 0.094 & 0.050 & 0.030 & 0.013 & 0.007 & 0.002 & 0.0006 & 0.170 \\
0.95 & 0.755 & 0.317 & 0.174 & 0.107 & 0.072 & 0.039 & 0.025 & 0.011 & 0.006 & 0.081 \\
\bottomrule
\end{tabular}
\end{center}
The floor rises with $\rho$ ($0.002$, $0.030$, $0.072$ at radius $4$), and at $\rho=0.95$ the radius-4 floor is close to the LMMSE floor itself: there, giving up the far sites costs about as much as giving up everything beyond second order. The Gaussian twin's windowed LMMSE reaches its full value by $r=3$ at $\rho=0.5$, so everything a wider window buys is non-Gaussian information. At $\rho=0.85$ the radius-4 floor peaks at $0.065$ around $t\approx0.7$--$1.0$ and is $0.030$ pooled---\emph{below} the window head's $0.118$ (all sites) and $0.075$ (interior) by a factor of two and a half to four, and close to EM--BP's $0.023$. Verdict: locality is \emph{not} the ceiling the manuscript and the report describe. A radius-4 estimator can in principle come within $0.007$ pooled of the exact-inference estimator; the window head's plateau is a property of that one-layer architecture and its training under the recipe, and the manuscript's sentence ``a radius-4 window cannot represent a message that has travelled the length of the chain'' is withdrawn (the message's information beyond radius 4 is worth $0.030$). What survives, sharpened: locality is a good prior at small $M$, and the plateau of the local architecture at large $M$ is an approximation limit that the exact floor now quantifies. The TCN (receptive radius 15) at $0.068$ and the global architectures at $0.060$ are likewise far above their floors ($0.0006$ at $r=16$; $0$ for the full chain), which says the same thing about all of them: at $M=2048$ under this recipe, every network is limited by approximation and optimisation, not by what its receptive field carries.

\subsection{The shape-information constant (E21, part ``constant'')}
\label{sec:scalar-const}
Scalar generalised Gaussian, untruncated quadrature: $I_t(\beta)e^{8t}/[\tfrac1{24}(\partial_\beta\kappa_4)^2]=1.0025, 1.0003, 1.0000$ at $t=1, 1.5, 2$ for $\beta=2$; $1.042, 1.006, 1.001$ for $\beta=3$; $0.960, 0.995, 0.999$ for $\beta=1.5$; the variance coordinate's $\tfrac12$ to five digits. The constant of eq.~(5) is correct. On the chain (E01) the outer-product estimate plateaus at $0.81$ ($\beta=2$) and $0.92$ ($\beta=3$) of the prediction from $t=3$ on, and diverges upward for $\beta\le1.5$ beyond $t\approx2.5$: the estimator's floor (a cancellation of $O(1)$ terms to $\alpha^4\approx10^{-5}$ against $10^{-5}$ grid accuracy) is reached before the asymptote for the kinked members and biases the smooth ones. The manuscript may say ``rate verified on the chain; constant verified in the scalar case and to $20\%$ on the chain''.

\subsection{The fourth-order correction (E16)}
\label{sec:e16}
Laplace chain, fraction of the squared closure gap closed by $\Sg+\nabla\epsilon/(1+\epsilon)$: $-11$ at $t=0.05$ (the expansion diverges; $\epsilon$ reaches $43$), $-1.4$ at $0.15$, $+0.30$ at $0.22$, $0.83$ at $0.32$, $0.92$ at $0.47$, $0.96$ at $0.68$, $0.985$ at $0.98$, $0.996$ at $1.4$, $1.000$ from $2.1$. Bimodal and uniform: the same shape, converging from $t\approx0.15$. Student-$t_5$: $0.13$ at $0.68$, $0.59$ at $1.4$, $0.82$ at $2.1$, then $0.39, 0.08, 0.01$ at $3, 3.5, 4$---on the $(801,8)$ grid, where the reference is truncation-dominated at those levels (Section~\ref{sec:e18}); a rerun at the wide domain is a one-line change and must be made before the number is used. Pooled over the twelve levels the corrected score is \emph{worse} than the twin ($0.47$ vs $0.155$ on Laplace) because the lowest levels dominate the pooled energy and the expansion diverges there: the fourth-order oracle is a per-level statement about what the gap is made of, not a competitor on the schedule. That is the right way to present it.

\subsection{Closed-form misspecification bias (E17)}
\label{sec:e17}
\begin{center}\small
{\small\setlength{\tabcolsep}{4pt}\begin{tabular}{@{}lrrrrrrr@{}}
\toprule
mechanism & $\varepsilon$ & $\rho^\ast$ & EM $\hat\rho$ & $b(\varepsilon)$ & lag-1 plug-in & EM--BP $512/2048$ & pred.\ $512/2048$ \\
\midrule
AR(2) & 0.1 & 0.867 & 0.855 & 0.060 & 0.052 & 0.056 / 0.055 & 0.063 / 0.061 \\
 & 0.4 & 0.916 & 0.888 & 0.248 & 0.216 & 0.214 / 0.214 & 0.248 / 0.248 \\
site-varying & 0.05 & 0.848 & 0.836 & 0.088 & 0.088 & 0.094 / 0.094 & 0.090 / 0.088 \\
 & 0.14 & 0.850 & 0.804 & 0.268 & 0.268 & 0.282 / 0.287 & 0.269 / 0.269 \\
long-range & 0.1 & 0.792 & 0.746 & 0.139 & 0.111 & 0.120 / 0.120 & 0.140 / 0.139 \\
 & 0.4 & 0.713 & 0.659 & 0.235 & 0.202 & 0.206 / 0.204 & 0.235 / 0.235 \\
rank-one & 0.5 & 0.886 & 0.884 & 0.028 & 0.022 & 0.030 / 0.027 & 0.034 / 0.030 \\
 & 1.0 & 0.937 & 0.932 & 0.065 & 0.048 & 0.055 / 0.052 & 0.068 / 0.066 \\
\midrule
\multicolumn{8}{@{}l}{small-$\varepsilon$ laws: AR(2) $b\approx0.60\,\varepsilon$; site-varying $b\approx1.75\,\varepsilon$;}\\
\multicolumn{8}{@{}l}{long-range $b\approx1.07\,\varepsilon^{0.8\to1.0}$; rank-one $b\approx0.15\,\beta^{2.00}$}\\
\bottomrule
\end{tabular}}
\end{center}
The KL projection under the one-view protocol reproduces EM--BP's $M=2048$ error within $0.003$ for site-varying $\rho$ and within $10$--$15\%$ for AR(2) and long-range, where the lag-one plug-in AR(1) (E08's \texttt{lmmse\_ar1} arm) is the closer match: the fitted $\hat\rho$ sits between the two projections because the $C=8$ kernel is not exactly Gaussian and the finite-sample fit under-shoots the KL optimum. Either way the asymptotic bias is a closed-form function of $\varepsilon$ and $L$; its order is $\varepsilon^1$ for the three mechanisms that perturb $O(L)$ directions of the covariance and $\varepsilon^2$ (in $\beta$) for the rank-one latent, which perturbs one. With $v=0.2$ from the well-specified cells, $\sqrt{b^2+v/M}$ predicts the two-size pattern of E08 at every strength.

\subsection{Oracle-free checkpoint selection (E20)}
Eight seeds at $M=32$, the size at which the choice of rung matters most (at $M=2048$ the frozen validation curves are monotone through the ladder, $0.0147\to0.0145$ over the last $400$ iterations, so every rule picks the top). Three seeds in: the oracle rule (validation risk against the true kernel) and the held-out one-view log-evidence rule select the \emph{same} rung in every case (rungs $20$, $40$, $40$), so the test errors are identical ($0.130$, $0.050$, $0.121$), against $0.232$, $0.187$, $0.158$ at the last rung of the ladder. The oracle is not what makes the selection work; the held-out evidence sees the same optimum. \textsc{[five seeds pending; each cell costs about forty minutes because the mixture M-step at $N_g=401$ is five seconds per iteration whatever $M$ is.]}

\subsection{Moment estimators (E19)}
\label{sec:e19}
Same chains and noise draws as E06, $M=512$, sixteen seeds.

\begin{center}\small
\begin{tabular}{@{}lrrrrrrrr@{}}
\toprule
$t$ & 0.05 & 0.1 & 0.2 & 0.4 & 0.7 & 1.0 & 1.4 & 2.0 \\
\midrule
RMSE $\hat\rho$, moments & 0.031 & 0.023 & 0.030 & 0.039 & 0.049 & 0.060 & 0.16 & 0.39 \\
RMSE $\hat\rho$, EM--BP & 0.005 & 0.005 & 0.006 & 0.008 & 0.011 & 0.009 & 0.03 & 0.34 \\
RMSE $\hat\kappa_4$, moments & 1.42 & 1.21 & 1.82 & 1.84 & 3.78 & 15 & 12700 & 1220 \\
RMSE $\hat\kappa_4$, EM--BP & 0.28 & 0.48 & 0.33 & 0.73 & 1.55 & 2.79 & 2.92 & 3.06 \\
\bottomrule
\end{tabular}
\end{center}
 Unbiased (bias $\le0.4$ up to $t=0.4$), never collapsing, $5$--$30\times$ the reference in dispersion at the informative levels, and with the asymptotic growth visible where the estimate is already useless: OLS log-slopes on $t\ge1$ of $1.83$ ($\hat\rho$, predicted $2$) and $3.73$ ($\hat\kappa_4$, predicted $4$). This is the high-noise competitor of the brief's Part~VII.6, and it answers cleanly: rate-correct, inefficient, and the inefficiency is the value of exact inference below the crossover (at high $t$ the MLE's score is itself the degree-4 Hermite statistic, so the two estimators become asymptotically equivalent---one sentence in the paper).

\subsection{The information curve at $\rho=0.5$ (E21, part ``curve'')}
\label{sec:e21}
Table~\ref{tab:crb}. Local slopes of $\log(1/I^{\mathrm{eff}}_\beta)$: $5.5, 5.7, 6.1, 6.5, 6.8, 6.9, 6.6$ on $[0.4,1.2]$; of $\log(1/I^{\mathrm{eff}}_\rho)$: $3.8$--$3.9$ throughout. E12b's levels $[0.6,1.0]$ sit where the shape slope is $6$--$7$. The corollary's exponents are the $t\to\infty$ limits; the experiment measures the curve. Amendment recorded in \texttt{experiments/registry/E12b.json} before any cell of that job ran.

\subsection{Interfaces and paired factors (E03 re-read)}
Table~\ref{tab:matched}. The ISM window head is the matched arm at every size; its paired factor at $M=2048$ is $7.9$ $[6.6,9.3]$; the fixed-corruption DSM arm, which sees the clean chain, is at $7.3$ $[6.2,8.7]$. Ratio-of-means versions: $7.2$ and $6.7$. Both should be given, with the estimand named as in Appendix~C.

\subsection{Compute accounting}
Single-thread timings on the laptop (under load; relative values are what matter), $L=32$:

\begin{center}\small
\begin{tabular}{@{}lrrr@{}}
\toprule
 & $N_g=201$ & $N_g=401$ & $N_g=801$ \\
\midrule
BP inference, one sequence, one level & 0.8 ms & 1.8 ms & 6.2 ms \\
E-step, one sequence & 1.0 ms & 2.5 ms & 9.1 ms \\
mixture M-step ($C=8$), per EM iteration & 1.5 s & 5.2 s & 22 s \\
\midrule
\bottomrule
\end{tabular}\\[6pt]
\begin{tabular}{@{}lrrrrr@{}}
\toprule
 & window & bi-GRU & U-Net & TCN & transformer \\
\midrule
forward pass, one sequence, one level & 0.11 ms & 0.20 ms & 0.22 ms & 0.33 ms & 0.41 ms \\
training, 32k steps (frozen \texttt{wall}) & 113 s & 525 s & 630 s & 641 s & 481 s \\
ISM training, 32k steps & 3{,}400 s$^\dagger$ & --- & --- & --- & 1{,}700 s$^\ddagger$ \\
\bottomrule
\end{tabular}\\[2pt]
{\footnotesize $^\dagger$ exact divergence; $^\ddagger$ Hutchinson estimator.}
\end{center}
So at the operating grid: inference over the twelve-level schedule costs EM--BP $\approx21$ ms per sequence against $1$--$5$ ms for a network (a factor $5$--$20$); an EM--BP fit at $M=2048$ with $\approx1000$ iterations costs $\approx(5+5)\,\mathrm{s}\times1000\approx3$ hours single-threaded against $2$--$11$ minutes for a network (a factor $\approx20$--$100$), with the mixture M-step, not the E-step, dominating below $M\approx2000$; memory is $O(N_g^2)$ per kernel and $O(ML N_g)$ per E-step chunk. The Pareto sentence for the paper: a factor $\approx8$ in error at equal data, bought with one structural assumption (priced by E17) and one to two orders of magnitude of compute. Scaling: linear in $L$ and $M$, quadratic in $N_g$, linear in $C$ for the M-step.

%% ------------------------------------------------------------------------------------
\section{New experiments, ranked}
\label{sec:plan}

Ranking is by scientific value per compute per referee risk removed. \hi{P0} = before submission; \hi{P1} = high value; \hi{P2} = optional. Costs are laptop core-hours (lc-h) or cluster node-days (nd). Seven of the items below are already registered and either done or running (E15--E21); the cluster queue holds E06c, E07b/c, E11b, E12b and E14.

\subsection*{P0}

\begin{enumerate}[label=\textbf{P0.\arabic*}]
\item \hi{E14, matched comparison across $\rho\in\{0.5,0.85,0.95\}\times L\in\{32,128\}$ (queued, cluster, $\approx2$~nd).} Hypothesis: EM--BP beats the best re-screened network at every cell at interface~C. Design as registered; eight seeds; estimand the paired factor. Expected under the ``property of the method'' reading: every interval excludes one and the factor is largest at $\rho=0.5$ (most fourth-order structure in the marginal, shortest posterior range) and smallest at $0.95$; under the ``peculiarity of the corner'' reading a cell near one or below. Supports or kills the headline. One addition to make \emph{now}, before it runs, at no cost: record the ISM objective's held-out curve so that a diverged ISM fit ($M=32$-style) can be recognised and reported as such rather than as a network's error.
\item \hi{E06c, extended ladder (running, cluster).} Resolves whether the $\hat\rho$ departure from the reference at $t=1.4$ and $2.0$ is budget. First cells (eleven of 144): the collapsed shape stays collapsed, but $\hat\rho$ moves substantially when the ladder is extended ($0.577\to0.800$ at $t=2$ for seed 0, rung $220\to3746$; $0.831\to0.844$ at $t=1.4$ for seed 1). Expected outcome: the $\hat\rho$ RMSE at $t=1.4$ and $2$ falls toward the reference and the ``leaves the bound'' clause is killed; the sentence becomes ``the correlation estimate follows the reference through $t=2$ once the optimiser is given the iterations the flat likelihood needs there''.
\item \hi{E12b with the amended analysis (queued, cluster).} Hypothesis now: $M^\ast(t)\propto1/I^{\mathrm{eff}}_t$ per coordinate, with local slopes $3.8$--$3.9$ ($\rho$) and $6.1$--$6.9$ (shape) on the swept levels (E21). Power: $1/I_\beta$ spans a factor $14$ over $[0.6,1.0]$, the $M$ ladder a factor $32$; with sixteen seeds and the exact enumerated bootstrap the separation of the two coordinates' slopes ($\approx2.7$ apart) is resolvable if $M^\ast$ is read at a threshold inside the identifiable window (CRB$(\kappa_4)<1$ needs $M\ge2048$ at $t=1.0$). The pre-registered thresholds should be checked against Table~\ref{tab:crb} before the data are read.
\item \hi{E07b/E07c, the family ladder at $A=12$ and $16$ (queued).} Expected: Student-$t$ rows move by the exceedance mechanism; at $A=16$ the expected number of sites beyond the domain is $0.04$ per replicate. Replace the flagged row; state the exceedance count per cell in the appendix.
\item \hi{E18 (done on the laptop, $\approx1$~lc-h).} Replaces Fig.~2c. Value: removes a false sentence and adds the heavy-tail proposition's evidence.
\item \hi{E15 (done, $\approx1$~lc-h).} Decides the locality sentence. Value: a referee who knows that a radius-$r$ posterior mean is computable will ask for exactly this.
\item \hi{E20 (done, $\approx5$~lc-h).} Oracle-free selection. Result: at $M=32$ the held-out evidence and the oracle pick the same rung in seven of eight seeds and neighbouring rungs in the eighth; mean test error $0.105$ under both against $0.208$ at the last rung (Section~6.8). The $M=512$ arm was dropped for cost. Kills the ``model selection uses the truth'' objection for EM--BP.
\end{enumerate}

\subsection*{P1}

\begin{enumerate}[label=\textbf{P1.\arabic*}]
\item \hi{Closure panel at the wide domain for all families to $t=5$ (E18 chain part; done).} Also gives the double-precision floor of the BP means (reached near $t=5$ for every family), which should be drawn as a floor in the figure.
\item \hi{E16 as a baseline row (done).} The fourth-order oracle score at each level: a ``moment-method'' estimator with oracle moments. Value: it explains \emph{what} the closure gap is made of at each level (fourth order above $t\approx0.3$, everything else below), and it is a natural bridge to the Balanov-type moment methods the related-work section invokes. Cost: none further. A version with \emph{estimated} $\kappa_4$ (from E19's moment estimator) would make it a real competitor: $\approx0.5$~lc-h; recommended.
\item \hi{Misspecification with the closed-form bias in the figure (E17; done).} Draw $b(\varepsilon)$ and $\sqrt{b^2+v/M}$ on Fig.~5 beside the measured EM--BP error; the break-even becomes a predicted intersection. Cost: figure work only.
\item \hi{Bulk-region companion for the locality story.} E03 already records bulk risks (window head $0.075$ at $M=2048$ on the bulk against $0.118$ all-site). Report both; the boundary fraction of the head's error is a fact worth one sentence.
\item \hi{Compute accounting (P1, $\approx0.2$~lc-h).} Inference $O(LN_g^2)$ per sequence per level; EM cost per iteration; network training wall-clock from the frozen \texttt{wall} columns; a Pareto table (data efficiency $\times$ compute $\times$ assumptions). Script prepared (\texttt{scratchpad/f\_compute\_cost.py}); run on an idle machine so the timings are clean.
\item \hi{Moment estimators in the learning figure (E19; done).} One extra pair of curves in Fig.~3: the unbiased, inefficient competitor that shows the $e^{2t}$, $e^{4t}$ growth without saturating (OLS slopes $1.83$ and $3.73$ on $t\ge1$) at $5$--$30\times$ the reference.
\end{enumerate}

\subsection*{P2}

\begin{enumerate}[label=\textbf{P2.\arabic*}]
\item \hi{A second innovation family for the matched comparison} (bimodal, where the fourth cumulant is negative): one $(\rho,L)$ cell, eight seeds, $\approx0.3$~nd. Tests whether the advantage is about heavy tails or about fourth-order structure of either sign.
\item \hi{A tree instead of a chain.} Not recommended: Theorem~2 is model-free, Theorem~3 is about $P_0$, and exact inference on a tree adds nothing to the principle; it adds a second laboratory and a page.
\item \hi{Lean: the projection operator of Proposition~1.} The report's decision~4. Not needed for the paper; the informal proof is complete.
\end{enumerate}

\begin{table}[h]\centering\small
\caption{E21: Cram\'er--Rao reference for the excess kurtosis at $\rho=0.5$ (Laplace member, one view per sequence), $\mathrm{d}\kappa_4/\mathrm{d}\beta/\sqrt{MI^{\mathrm{eff}}_\beta}$; the coordinate's range is $3$ and a bound above $\approx1$ leaves no usable window. Local slopes of $\log(1/I^{\mathrm{eff}})$ between consecutive rows in the last two columns. Source: \texttt{outputs/E21/information.csv}.}
\label{tab:crb}
\begin{tabular}{@{}rrrrrrrr@{}}
\toprule
$t$ & $M=256$ & $512$ & $2048$ & $8192$ & $16384$ & slope ($\kappa_4$) & slope ($\rho$) \\
\midrule
0.6 & 1.24 & 0.88 & 0.44 & 0.22 & 0.16 & 6.1 & 3.8 \\
0.7 & 1.68 & 1.19 & 0.59 & 0.30 & 0.21 & 6.5 & 3.8 \\
0.8 & 2.32 & 1.64 & 0.82 & 0.41 & 0.29 & 6.8 & 3.8 \\
0.9 & 3.26 & 2.31 & 1.15 & 0.58 & 0.41 & 6.9 & 3.8 \\
1.0 & 4.59 & 3.25 & 1.62 & 0.81 & 0.57 & 6.6 & 3.9 \\
1.2 & 8.90 & 6.29 & 3.15 & 1.57 & 1.11 & 5.5$^\ast$ & 3.9 \\
1.4 & 15.5 & 11.0 & 5.48 & 2.74 & 1.94 & --- & 3.9 \\
\bottomrule
\end{tabular}\\[2pt]
{\footnotesize $^\ast$ the outer-product estimate of $I_\beta$ reaches its quadrature floor beyond $t\approx1.2$ at $N_g=401$ for the kinked Laplace member; slopes there are not informative.}
\end{table}

%% ------------------------------------------------------------------------------------
\section{Recommended final claims}
\label{sec:claims}

\subsection{The strongest defensible version of each headline claim}

\begin{enumerate}
\item \hi{Laboratory.} ``For a first-order Markov chain with continuous non-Gaussian transitions under coordinatewise OU noise, the posterior factor graph is a tree and the diffusion score is computed exactly by one forward--backward sweep of functional messages; we carry the messages on a grid whose truncation and quadrature errors are measured, not assumed.''
\item \hi{Proposition 1.} ``Gaussian message closure computes the LMMSE estimator of the covariance twin for every innovation law with the same first two moments. Consequently every linear estimator, and every Gaussian-innovation kernel fitted by (quasi-)maximum likelihood, is bounded below by the LMMSE floor of the true covariance.''
\item \hi{Theorem 2.} ``Under the channel every Hermite moment of degree $m$ is attenuated exactly by $e^{-mt}$, assuming only finite $m$-th moments.'' Followed by the dictionary: amplitude $e^{-mt}$, information $e^{-2mt}$, samples $e^{2mt}$.
\item \hi{Theorem 3 (restated).} ``If $P_0$ has an exponential moment, then as $t\to\infty$ at fixed $L$: $\E\|S-\Sg\|^2=\tfrac16\Frob{\kappa_4}^2e^{-8t}(1+O(e^{-2t}))$; $I_t(\theta)=\tfrac12\Frob{\partial_\theta\Sigma_0}^2e^{-4t}(1+O(e^{-2t}))$ for a covariance coordinate; $I_t(\theta)=\tfrac1{24}\Frob{\partial_\theta\kappa_4}^2e^{-8t}(1+O(e^{-2t}))$ for a shape coordinate; in general $e^{-2mt}$ times the Hermite norm of the derivative of the degree-$m$ moment tensor (Theorem~\ref{thm:general}). For the Gaussian chain the correction is explicitly $O(e^{-2t}\lmax)$ and $\tstar=\tfrac12\log\lmax$ is the exact scale.''
\item \hi{Verification.} ``The Gaussian-chain rate and constant are reproduced exactly; the score-gap constant to $1$--$6\%$ on the chain for every family on a domain wide enough for its tail, including Student-$t_5$ (to $1\%$ at $t=4$); the shape rate to $-8.01$ and its constant to $0.1\%$ in the scalar case and to $10$--$20\%$ on the chain, where the outer-product estimator's floor limits the check.''
\item \hi{Heavy tails.} ``For a regularly varying innovation of index $\nu$ the expected gap acquires an $e^{-\nu t}$ term carried by events of probability $\sim e^{-\nu t}$; on a finite sample the Hermite law is what is seen. Theorem~3's hypothesis excludes such laws; the E01 departure at $t\approx3$ was truncation.''
\item \hi{Corollary (sample complexity).} ``In the local asymptotic minimax sense, at fixed level $t$, estimating a degree-$m$ coordinate to fixed local precision needs $\Theta(e^{2mt})$ sequences, and maximum likelihood attains the constant; a shape coordinate needs $e^{4t}$ times more than a covariance coordinate.''
\item \hi{Fourth-order correction (new).} ``Adding the degree-4 Hermite term back to the twin closes $96\%$ of the closure gap at $t=0.68$, $99.6\%$ at $1.4$ and all of it by $2$ on the Laplace chain; the correction diverges below $t\approx0.2$; on Student-$t$ it closes $82\%$ at $2.1$ and $8\%$ at $3.5$.''
\item \hi{Learning through noise.} ``Fitting the transition law from $M=512$ one-view sequences at level $t$: for $t\le0.4$ the fitted shape is unimodal with dispersion within $20\%$ of the one-parameter Cram\'er--Rao reference; at $t=0.7$ a quarter of the fits collapse to a Gaussian innovation and the rest are shrunk; at $t=1$ all but one collapse, one level below the level at which the reference bound crosses the coordinate's range. The correlation is recovered to $0.01$ through $t=1$ with dispersion at the reference. Moment estimators of the same coordinates from the same data are unbiased at $5$--$30\times$ the reference and show the $e^{2t}$, $e^{4t}$ growth.''
\item \hi{Capacity.} As now: ``$C=1$ sits on the floor at every size and across the family and the regime sweep ($0.01$ worst); $C=8$ is $7.9\times$ below it at $M=2048$; the previous conclusion was a stopping-rule artefact, diagnosed and controlled.''
\item \hi{Structured vs neural, standard interface.} ``Under one recipe the structured estimator is ahead of six architectures at every size, by a paired $2.8\times$ $[2.4,3.4]$ against the best at $M=2048$; the networks with the most headroom were still improving at 32k steps, so the factor is an upper bound.''
\item \hi{Structured vs neural, matched interface.} ``At the estimator's own interface---one noised view per sequence, no clean chain, checkpoint selected without the true score---the best network reaches $0.165$ against $0.023$, a paired $7.9\times$ $[6.6,9.3]$ in $16/16$ seeds; even given the clean sequence paired with its single corruption it reaches $0.155$. The advantage at equal information is larger than at the standard interface, not smaller.''
\item \hi{Locality.} Pending E15's pooled number at $\rho=0.85$. Two candidate sentences: (a) if the radius-4 floor is $\ge0.06$: ``the window head's plateau is within a factor two of the exact radius-4 floor; locality is a ceiling''; (b) if $\le0.03$: ``the exact radius-4 estimator is far below the window head; the plateau is a limit of that architecture and its training, and the information a radius-4 window carries is mostly non-Gaussian and mostly unexploited''.
\item \hi{Length.} As now.
\item \hi{Generation.} ``The fitted kernel generates within $0.09$ nats $[0.00,0.19]$ of the exact score and with excess kurtosis $1.80$ against $1.94$; the covariance twin generates a Gaussian chain ($0.00$) at $+1.70$ nats; sliced $W_2$ cannot tell them apart.''
\item \hi{Misspecification.} ``The asymptotic bias of the chain estimator under a Gaussian departure is a closed-form KL projection that reproduces the $M=2048$ error to $\pm0.003$; it is $\Theta(\varepsilon)$ for AR(2), site-varying $\rho$ and long-range coupling and $\Theta(\varepsilon^2)$ (and $O(1/L)$) for a rank-one latent, which is why the last never reverses the ordering in the swept range and the first two reverse it at $\rho_2\in(0.3,0.4)$ and $\delta\in(0.10,0.14)$.''
\end{enumerate}

\subsection{Claims that must not appear}

\begin{itemize}
\item ``The estimator attains the Cram\'er--Rao bound'' / ``$1.02\times$ the bound''.
\item ``Information-matched'' for the fixed-corruption DSM arm; ``$7.3\times$ at matched information''.
\item ``The experiment confirms the $e^{4t}$ separation'' / ``the $e^{4t}$ separation \ldots measured''.
\item ``Student-$t_5$ departs as the moment condition predicts'' (at $t\approx3$).
\item ``Essentially all of its belief mass outside the domain''.
\item ``Requires $M\propto e^{8t}$ samples'' without ``local asymptotic minimax''.
\item ``Finite eighth moments'' as the hypothesis of Theorem~3 with the present proof.
\item ``What is not in that literature is the score-level statement for a Markov (dependent) rather than i.i.d.\ latent structure'' (the low-SNR expansions are for arbitrary inputs).
\item ``A radius-4 window cannot represent a message that has travelled the length of the chain'' unless E15 supports it.
\item ``The generator refuses cells whose narrowest component is under two grid cells'' (it flags them).
\item ``Sixteen paired seeds are the inferential unit'' as a blanket statement (eight in E13/E14).
\item ``We are the first''---anywhere.
\end{itemize}

\subsection{Novelty, stated precisely}

\emph{Established before this paper:} Mehler's formula and OU spectral theory; cumulant attenuation $e^{-nt}$ under the forward process (Aarts et al.\ 2025); low-SNR expansions of mutual information and MMSE with fourth-order statistics at second order (Prelov--Verd\'u 2004; Guo--Shamai--Verd\'u 2005, 2011), valid for arbitrary input laws; the SNR$^k$ information law for directions first identifiable at moment order $k$ in latent-variable models (Katsevich--Bandeira 2023; Perry et al.\ 2019; Bandeira et al.\ 2023; Balanov et al.\ 2026); consistency and identifiability of MLE for general-state HMMs (Douc et al.\ 2011; Gassiat et al.\ 2020); EM via Fisher's identity (Capp\'e et al.\ 2005); Gaussian-dominance observations for learned scores (Wang--Vastola 2024; Li et al.\ 2024).

\emph{Specialisations of known theory here:} eqs.~(4)--(5) as the diffusion-time form of the SNR$^k$ law; the LAN corollary.

\emph{New in this paper:} (i) an exactly solvable, correlated, non-Gaussian diffusion laboratory with a validated instrument; (ii) Proposition~1 (closure = LMMSE of the twin) and its use as an apparatus check across fifteen regimes; (iii) the score-level statement of Theorem~3 with its constant $\tfrac16\Frob{\kappa_4}^2$, the general-order constant (Hermite norm), the exponential-moment hypothesis with the heavy-tailed regime beyond it, and the crossover as an exact scale for the Gaussian chain; (iv) exact verification of rates and constants with the true score; (v) transition learning by exact-inference EM through the channel, with the collapse level measured against the information curve, the moment competitor, and the oracle-free selection control; (vi) the controlled structured-vs-neural comparison at three declared interfaces with a matched interface that gives the larger factor; (vii) the fourth-order oracle correction as the moment-method reading of the closure gap; (viii) the closed-form misspecification bias and its order in $\varepsilon$. The contribution is the synthesis; no sentence should claim more.

%% ------------------------------------------------------------------------------------
\section{Recommended ICML narrative}
\label{sec:narrative}

\subsection{What the paper is about}

The four candidate stories of the brief (A: noise erases information order by order; B: exact structured inference measures it; C: structure converts surviving information into sample efficiency; D: the gain lasts while the structure is approximately right) are one story read at four distances, and the evidence now supports the hierarchy the brief suspected, with one change of emphasis. The load-bearing new fact is not the $7.9\times$; it is that every learning result in the paper is \emph{predicted} by the information curve: the collapse level of the shape estimate (E06 vs the CRB crossing), the misspecification break-even (E17 vs the network's error), the closure gap and what it is made of (E16), the plateau of a local architecture (E15, pending), and the sample-complexity slopes (E21 vs E12b). The paper should be organised so that a reader sees the information theory make a prediction and the exact laboratory test it, five times.

\subsection{The arc}

\begin{enumerate}
\item \hi{Question} (\S1). What information about structured data survives diffusion noise, and what does knowing the structure buy? Answered in one sentence each for theory (order by order, at explicit rates and constants), inference (exact structured inference recovers what is left, and its error is at the information curve wherever the coordinate is identifiable), and its limit (a wrong structure is a bias no data removes, of an order in $\varepsilon$ we compute).
\item \hi{Laboratory} (\S2). The chain, the channel, the tree, the exact score, the covariance twin, Proposition~1 with its two corollaries (every linear estimator is bounded by the floor; a Gaussian kernel fitted by likelihood lands on it). One figure: the factor graph.
\item \hi{Theory} (\S3). Theorem~2 with the dictionary. Theorem~3 under Assumption~(E), with the general-order theorem and the LAN corollary; the crossover as an exact scale for the Gaussian chain. Verification: Figure~2 (three panels, wide-domain closure). The heavy-tailed regime in two sentences with Proposition~\ref{prop:heavy}. The fourth-order oracle (E16) in one sentence: ``adding the term back closes the gap''---the theorem's object is the gap's content.
\item \hi{Learning through the channel} (\S4). EM around the exact recursion; Fisher's identity. Figure~3 redesigned: the estimator's dispersion and its collapse fraction against the information reference per level, with the moment estimator as the rate-correct, inefficient competitor. The claim: the estimator is at the reference while the coordinate is identifiable and collapses where the reference says it must; the exponents are the curve's, and the curve is exact.
\item \hi{What explicit structure buys} (\S5). Capacity ($C=1$ on the floor everywhere, $C\ge2$ far below; the retired reading with its diagnosis in one paragraph); the family $V$; the three interfaces against six architectures with the matched (ISM) factor as the headline and the regime sweep E14 as its robustness; locality against the exact radius floor (whatever E15 says, it is a cleaner statement than a plateau); length; generation in three numbers.
\item \hi{Where it stops} (\S5, last paragraph, or \S6). Misspecification with the closed-form bias drawn on the figure: the break-even is a predicted intersection; the order in $\varepsilon$ explains which mechanisms bite.
\item \hi{Related work} (\S6), two thirds of a column, with the novelty stated as the synthesis of Section~\ref{sec:claims}.
\item \hi{Limitations} (\S7): one dimension, linear autoregression, $O(LN_g^2)$ inference (with the Pareto sentence: data efficiency against compute and assumptions), fixed $L$ in the theory, eight seeds in the sweeps, and the heavy-tailed case outside Theorem~3.
\end{enumerate}

%% ------------------------------------------------------------------------------------
\section{Proposed paper structure}
\label{sec:structure}
\subsection{Section by section, with a page budget}

\begin{tabular}{@{}lp{9.6cm}r@{}}
\toprule
section & content & pages \\
\midrule
1 Introduction & question; three answers; contributions (six items) & 0.9 \\
2 Laboratory & model, channel, tree, exact score, twin, Prop.~1 + corollaries, Fig.~1 & 0.9 \\
3 Theory & Thm~2 + dictionary; Thm~3 (E); Thm~\ref{thm:general}; Cor.~\ref{cor:lam}; crossover; heavy tails; verification Fig.~2; the fourth-order oracle & 2.0 \\
4 Learning through the channel & EM, Fisher's identity; Fig.~3 (dispersion, collapse, reference, moments) & 1.0 \\
5 What structure buys & capacity (Table~1, Fig.~4 family); interfaces (Table~2, matched headline); locality vs floor; length (Table~3); generation (3 numbers); misspecification (Fig.~5 with $b(\varepsilon)$) & 2.2 \\
6 Related work & as \S8.3 of this report & 0.5 \\
7 Limitations and conclusion & & 0.4 \\
\bottomrule
\end{tabular}

\noindent Appendices: A proofs (with the corrected tail argument and the general-order proof); B numerics (envelope by exceedance count; the uniform kernel's $O(h)$ note; the double-precision floor of the means); C protocol (three interfaces named; selection rules named; seeds per experiment); D references and oracles; E the retired capacity reading; F E14 regime table when it lands; G the closed-form bias table; H the fourth-order oracle per level. Drop the archival aggregation appendices (D/G/H of the current draft), as the plan already says.

%% ------------------------------------------------------------------------------------
\section{Figures and tables redesigned}
\label{sec:figures}


Each visual must prove one claim, mark $\tstar$ and floors where they apply, and separate pre-asymptotic from asymptotic ranges.

\begin{itemize}
\item \hi{Fig.~1 (mechanism).} As now.
\item \hi{Fig.~2 (order-by-order theory), three panels.} (a) Gaussian-chain information scaled by its asymptote, five $\rho$, $\tstar$ ticks---as now. (b) Efficient information about $\rho$ and $\beta$ with $e^{-4t}$, $e^{-8t}$ guides \emph{and} the predicted constants drawn as horizontal asymptotes on a scaled axis, so the reader sees the constant reached at $\beta=1.5$ and missed by $20\%$ at $\beta=2$ with the floor marked. (c) Closure gap divided by the predicted constant for Laplace, bimodal, uniform, Student-$t_{5,7,9}$ from E18 at $(3201,32)$, with the double-precision floor of the means as a shaded region above $t\approx4.5$; the $(801,8)$ Student-$t$ curve as a thin grey line labelled ``truncated domain''. Proves: rates and constants; the truncation artefact; the heavy-tail family reaching the constant.
\item \hi{Fig.~3 (learning through noise), two panels.} (a) $\hat\rho$: dispersion (not RMSE) per level with the CRB reference and the moment estimator; (b) $\hat\kappa_4$: the fits themselves as points (so the bimodality is visible), the non-collapsed cluster's dispersion, the collapse fraction printed per level, the reference bound, the coordinate's range line, and the moment estimator. Proves: at the reference while identifiable; collapse where the reference crosses the range; rates in the competitor.
\item \hi{Fig.~4 (capacity), one panel.} The family $V$ as now (it is the best figure), with the regime sweep as an inset table or a second small panel: paired gain vs $\rho$ at $L=64$.
\item \hi{Fig.~5 (structure vs networks), two panels.} (a) Error vs $M$ for the six architectures at interface~A and EM--BP, with the LMMSE floor; (b) the matched interface: ISM window and transformer, the fixed-corruption arm dashed, EM--BP, at four sizes; the E14 regime table beside it when it lands. Proves: the factor at equal information, and that it is larger than at the standard interface.
\item \hi{Fig.~6 (locality and misspecification), two panels.} (a) The exact radius-$r$ floor vs $r$ at $\rho=0.85$ (pooled and at two levels) with the window head, TCN (radius 15) and EM--BP as horizontal marks: locality is read against its floor. (b) Misspecification: the four mechanisms as now, with the closed-form bias $b(\varepsilon)$ and the finite-$M$ prediction drawn through the EM--BP points; the break-even as the intersection with the network's curve.
\item \hi{Table~1 (capacity)} as now with the flag policy stated in the caption. \hi{Table~2 (interfaces)}: rows = arms, columns = $M$, grouped by interface A/B/C, EM--BP last, LMMSE floor first. \hi{Table~3 (length)} as now, adding the bulk-region column for the window head. \hi{Table~4 (generation)}: cut to $M=2048$ and three columns ($\Delta\ell$, $\hat\kappa_4$, $W_2$); the $M=128$ block to the appendix.
\item \hi{Remove:} the phase-ladder figure from the body (appendix), the aggregation tables (drop), the archival comparison appendix (drop).
\end{itemize}

%% ------------------------------------------------------------------------------------
\section{Reviewer simulation}
\label{sec:reviews}

\subsection{The strongest rejection review (of the paper as submitted on 18 September)}

\emph{Summary.} The paper studies a Markov chain with non-Gaussian innovations under OU noise, proves Hermite attenuation and an order-by-order information decay, learns the transition by EM, and compares with neural denoisers. The laboratory is neat and the experiments are carefully run, but the theoretical contribution is a specialisation of known low-SNR results with a proof that does not hold under its stated hypothesis, and the headline empirical claims are overstated or mislabelled.

\begin{enumerate}
\item \emph{Novelty.} The $e^{-2mt}$ information law is the SNR$^k$ law of Katsevich--Bandeira and Balanov et al.\ in the diffusion clock; cumulant attenuation is in Aarts et al.; the second-order MMSE expansions of Prelov--Verd\'u and Guo et al.\ hold for arbitrary inputs, dependent ones included, so ``for a Markov rather than i.i.d.\ latent structure'' is not a novelty. What remains is a constant and a synthetic verification. \hi{Valid in part.} Fix: claim the synthesis (Section~\ref{sec:claims}); add what is new beyond the constant---the general-order theorem with its Hermite-norm constant, the heavy-tailed regime, the fourth-order oracle, the closed-form bias---each of which is a statement the cited papers do not make.
\item \emph{Theorem 3's proof.} Real-analyticity of $\Lambda$ at the origin and $P_t(|\alpha x|>1)\le e^{-c/\alpha^2}$ are false under finite eighth moments. \hi{Valid.} Fixed by Assumption~(E) (Section~\ref{sec:thm3}).
\item \emph{Sample complexity.} Corollary~5 turns a Fisher information into ``requires''. \hi{Valid.} Fixed by Corollary~\ref{cor:lam}.
\item \emph{Narrow setting.} One dimension, one chain length in the theory verification, one family in most experiments. \hi{Partly valid.} Answer: the theorems are for any $P_0$ with (E); the regime sweep covers five $\rho$ and three $L$; E14 covers three $\rho$ and two $L$ for the network comparison. Say it as a limitation, not a defence.
\item \emph{Baseline fairness.} The networks train for 32k steps and the two largest are still improving at the last checkpoint in $88\%$ of cells; the EM ladder was extended when it bound. \hi{Valid.} Fix: state the at-cap rates for the networks and call the standard-interface factor an upper bound; the matched arm is not affected (window head never at cap).
\item \emph{Information mismatch.} ``Matched'' arm sees the clean chain. \hi{Valid.} Fixed by making ISM the headline; the number grows.
\item \emph{Numerical approximation.} The reference is a grid; the Student-$t$ cell is flagged; the closure panel departs. \hi{Valid in part.} Fixed by E18 and the exceedance diagnostic; the frozen-configuration validation stands.
\item \emph{Seeds.} Eight in the sweeps. \hi{Valid as stated.} Fix: say so; the intervals are paired and exclude zero at eight.
\item \emph{Selection bias.} Every arm but one selects on the true score. \hi{Valid.} Fixed by E20 (and by making the oracle-free ISM arm the headline).
\item \emph{Computational impracticality.} $O(LN_g^2)$ per sequence. \hi{Valid and admitted.} Fix: the Pareto sentence with numbers (P1.5).
\item \emph{One parameter point.} The $7.3\times$ at one $(\rho,L)$. \hi{Valid until E14.}
\item \emph{Estimation vs optimisation.} The efficiency claim mixes EM's early stopping with the channel. \hi{Valid.} Fixed by the per-level table.
\item \emph{Connection to modern diffusion models.} The chain is far from images. \hi{Valid; not fixable by experiment.} Answer in the conclusion: the theorem is the general statement; the laboratory is what makes it testable to the constant; the Gaussian-dominance observations of Wang--Vastola and Li et al.\ are its images-scale shadow.
\end{enumerate}
\emph{Recommendation as submitted: reject; resubmit with the hypothesis fixed, the interfaces named, and the efficiency claim withdrawn.}

\subsection{The strongest acceptance review (of the paper this report recommends)}

\emph{Summary.} An exactly solvable, correlated, non-Gaussian diffusion laboratory; a clean theorem (Hermite attenuation) and a sharp asymptotic (order-by-order information decay with explicit constants, a general-order form, and a crossover that is exact for the Gaussian chain), verified with the exact score to the constant on six innovation laws; a learning experiment in which an exact-inference EM is placed against the information curve and collapses exactly where the curve says the coordinate stops being identifiable; a controlled comparison against six sequence architectures at three declared information interfaces, with the matched interface giving the larger factor and an oracle-free control; a closed-form misspecification bias that reproduces the measured errors and explains which departures bite at which order. The retired claims are recorded with their diagnoses.

\emph{Strengths.} (1) The only setting I know in which a non-Gaussian dependent score is available exactly and the Gaussian part of the score is identified with a named estimator (Proposition~1), turning ``learned scores are dominated by their Gaussian part'' into a rate with a constant. (2) The information theory predicts the experiments five times (collapse level, closure content, break-even, floor, slopes) rather than being fitted to them. (3) The experimental discipline: hashed bundles, paired seeds, registered analyses, an amendment recorded before data, a control for every confound a referee could name. (4) Honesty about what was retired and why.

\emph{Weaknesses.} One dimension; $O(LN_g^2)$ inference; eight seeds in the sweeps; the heavy-tailed case only numerically supported; the sample-complexity exponents measured through the curve, not through the asymptote.

\emph{Recommendation: accept.}

\subsection{Remaining risks after the fixes}

\begin{itemize}
\item E14 finds a cell where the network ties or wins (e.g.\ $\rho=0.95$, $L=128$, where the marginal is nearly Gaussian). Then the headline becomes ``the advantage tracks the fourth-order structure of the marginal'', which is the theory's prediction and is not a loss.
\item E15's floor at $\rho=0.85$ turned out far below the window head ($0.030$ against $0.118$), so the locality paragraph has changed from a ceiling to an architecture limit; a referee who checks will find the floor computed. The residual risk is the opposite one: a referee may ask why no local architecture reaches its floor, which is a question about training recipes this paper does not answer.
\item E12b's slopes come out at the curve's values ($\approx6.5$, $3.8$) and a referee reads them as ``not 8, not 4''. The amendment, recorded before data, and Table~\ref{tab:crb} are the answer.
\item The heavy-tailed proposition is challenged as unproved. It is labelled as numerically supported; the scalar quadrature (E18 scalar part) is the evidence; the paper should not promote it to a theorem.
\item Page budget: with the audit's additions the body ends on page 12 against a limit of 8 (the second row of Fig.~5 cost half a page; the locality figure sits in the appendix and costs none). The cut order in RESEARCH\_PLAN.md (phase figure and the $M=128$ generation block to the appendix, related work to two thirds of a column, the setup displays to one line) recovers about two pages; the rest must come from \S5's prose, which now carries three paragraphs of caveats that can each lose a sentence once E14 fixes the framing, and from the choice between the two rows of Fig.~5 if the cut is not otherwise found. This is editing, not science, and was deliberately left for after the cluster runs land.
\end{itemize}

%% ------------------------------------------------------------------------------------
\section{Final submission checklist}
\label{sec:checklist}

\begin{enumerate}[label=$\square$, leftmargin=1.6em]
\item \textsc{done} Theorem~3 restated under Assumption~(E); Appendix~A.3 tail paragraph rewritten (exponential tail from the exponential moment; tower property for the conditional second moment); ``finite eighth moments'' removed from the statement.
\item \textsc{done} Theorem~\ref{thm:general} (general order) in the body with its constant; Remark~4 removed.
\item \textsc{done} Corollary~\ref{cor:lam} replaces Corollary~5; ``requires'' removed.
\item \textsc{done} The dictionary sentence (amplitude / information / samples) in \S3 and the abstract rewritten accordingly; \texttt{[PENDING]} removed.
\item \textsc{done} Proposition~1's two corollaries (linear estimators; Gaussian QMLE consistency) added in one sentence each.
\item \textsc{done} Fig.~2c redrawn from E18 at $(3201,32)$ with the floor shaded; caption: Student-$t_5$ reaches the constant; the heavy-tailed regime in one sentence with Proposition~\ref{prop:heavy}; the ``departs as the moment condition predicts'' sentence removed.
\item \textsc{done} ``Verifies the rates and the constants'' replaced by the sentence that names which constant is verified where (E21 scalar; chain to $10$--$20\%$).
\item \textsc{done; Fig.~3 unchanged, already shows the individual fits} \S4: the efficiency paragraph replaced by the per-level statement (unimodal and at the reference for $t\le0.4$; collapse fraction $4/16$ at $0.7$, $15/16$ at $1.0$; reference crossing at $0.98$); Fig.~3 redrawn with dispersion, points, collapse fractions and the moment estimator (E19).
\item \textsc{done} \S1 and contributions: ``confirms'' / ``measured'' for the $e^{4t}$ separation removed; ``rates by computation, ordering by measurement'' in.
\item \textsc{done} \S5 matched interface: ISM arm as the headline ($7.9\times$ $[6.6,9.3]$), fixed-corruption arm as ``even with the clean target''; interfaces A/B/C named in the protocol; the $M=32$ ISM divergence noted; \texttt{make\_baseline\_numbers.py} changed to emit the ISM arm as the matched headline by definition.
\item \textsc{done} \S5 standard interface: at-cap rates of the networks stated; ``upper bound'' clause added.
\item \textsc{done; Fig.~6a not yet drawn} Locality paragraph rewritten against E15's floor (either sentence of Section~\ref{sec:claims}, item 13); bulk-region numbers cited; Fig.~6a added.
\item \textsc{done; $b(\varepsilon)$ not yet drawn on Fig.~5} E08 window head named (radius 2, width 64, 16k steps); the rank-one explanation replaced by the two reasons of Proposition~\ref{prop:bias}; Fig.~5 with $b(\varepsilon)$ and $\sqrt{b^2+v/M}$ drawn (E17).
\item \textsc{done; E07 row pending the cluster} App.~B envelope paragraph: exceedance count and mechanism; ``essentially all'' removed; the uniform kernel's $O(h)$ note; E07 Student-$t$ row replaced when E07b/c land, with the exceedance count per cell.
\item \textsc{done; E20 at $M=32$ pending} App.~C: flag policy for under-resolved fits stated correctly; seeds per experiment; selection rule per arm; E20 result cited.
\item \textsc{done} \verb|\cref{thm:gap}| fixed.
\item \textsc{done} Related work: novelty paragraph replaced by the synthesis statement; ``dependent rather than i.i.d.'' removed.
\item \textsc{done} Limitations: fixed $L$ in the theory; heavy tails; eight seeds; $O(LN_g^2)$ with the Pareto numbers.
\item \textsc{open: cluster} E14 folded in when it lands (Table or Fig.~5b); the framing sentence chosen after it, per the report's decision~7.
\item \textsc{open: cluster; first cells say budget} E06c folded in: the ``leaves the bound'' clause resolved.
\item \textsc{open: cluster} E12b read against Table~\ref{tab:crb}; the amendment cited.
\item \textsc{generators done; \texttt{make check} to be run after E15/E20 freeze} New generators for E15--E21 numbers (\texttt{tools/make\_audit\_numbers.py}) so no typed decimal enters the body; \texttt{make check} green; \texttt{audit\_claims.py} extended to the new macros; outputs frozen with \texttt{tools/freeze.py}.
\item \textsc{done} CLAIM\_EVIDENCE\_MATRIX rows E11, E16 rewritten; NEGATIVE\_RESULTS entries for the efficiency claim, the matched-interface label, the Student-$t$ departure and the belief-mass sentence; SCIENCE\_AUDIT entry for this audit.
\item \textsc{open: the body currently ends on page 11, three pages over; the recorded cut order recovers about two, the third must come from \S5's prose} Page budget $\le8$ body pages after the cuts; every figure rendered and looked at; no Overfull box from a table or figure.
\end{enumerate}


\end{document}
